\documentclass[11pt,a4paper]{article}
\usepackage[margin=2.5cm]{geometry}
\usepackage{amsmath,amssymb,amsthm}
\usepackage{graphicx}
\graphicspath{{./}{../github/fsri-theory/figures/}}
\usepackage{hyperref}
\usepackage{booktabs}
\usepackage{authblk}
\usepackage{abstract}
\usepackage{titlesec}
\usepackage[numbers]{natbib}
\usepackage{xcolor}
\usepackage{mdframed}

\titleformat{\section}{\large\bfseries}{\thesection.}{0.5em}{}
\titleformat{\subsection}{\normalsize\bfseries}{\thesubsection.}{0.5em}{}

\newtheorem{theorem}{Theorem}
\newtheorem{proposition}{Proposition}
\newtheorem{corollary}{Corollary}
\newtheorem{lemma}{Lemma}
\newtheorem{definition}{Definition}
\newtheorem{remark}{Remark}

\hypersetup{colorlinks=true,linkcolor=blue,citecolor=blue,urlcolor=blue}

% ============================================================
\title{\textbf{The Functional Symbiotic Resilience Index:\\
Topological Entropy, Wasserstein Curvature Bounds,\\
and Non-Equilibrium Thermodynamics of Complex Networks}}

\author[1]{Alberto Acedo}
\affil[1]{Biome Makers Inc. \quad \texttt{acedo@biomemakers.com}}
\date{August 2026}
% ============================================================

\begin{document}
\maketitle

\begin{abstract}
We present a mathematical framework unifying network topology,
non-equilibrium thermodynamics, and resilience theory through the
\textbf{Functional Symbiotic Resilience Index} (FSRI, $\Omega$).
The central object is the topological pressure functional
$\mathcal{T}(W) = \operatorname{Tr}(WW^\top)^3$, which we show to be
(i) proportional to the global clustering coefficient of the induced
weighted graph, (ii) the generator of a gradient flow whose fixed points
correspond to Prigoginian minimum-entropy-production states, and
(iii) a Gini-measurable concentration index whose gradient $3(WW^\top)^2$
identifies structural monopolies, that is, nodes or weights that accumulate
disproportionate topological influence.

We further establish where the framework's information actually resides.
The global functional obeys the exact identity
$\operatorname{Tr}(A^3) = \sum_i \lambda_i^3$: it \emph{is} the third
spectral moment, and is therefore a member of the same family as the
Absorption Ratio, effective rank and the Vendi score rather than an
alternative to it. By contrast, the node-wise decomposition
$(A^3)_{ii} = \sum_k \lambda_k^3 (v_k)_i^2$ depends on eigenvectors and is
\emph{not} determined by the spectrum: we give an explicit cospectral pair
with identical $\operatorname{Tr}(A^3)$ but distinct
$\operatorname{diag}(A^3)$. This locates the non-redundant content of the
framework in the node-wise field rather than in the global scalar, and
predicts the empirical pattern observed in applications. We also show that
on \emph{signed} matrices the same functional measures Heider structural
balance, information inaccessible to any index computed on $|A|$.

We validate these properties through exact analytical results and
computational experiments on synthetic networks (Barabási-Albert,
Watts-Strogatz, Erdős-Rényi, ring lattice), demonstrating that the
Gini coefficient of the topological pressure distribution distinguishes
monopoly networks ($G \approx 0.986$) from uniform networks
($G \approx 0.272$) by a factor of 3.6×.

The framework has an immediate computational application: the Omega-S
regularizer~\cite{acedo2026omegas}, which minimises
$\operatorname{Tr}(WW^\top)^3$ during neural network training via
Hutchinson stochastic trace estimation. In the applied work that penalty is
combined with group lasso for structured pruning, where the compression
obtained is attributable to the group-lasso term and the marginal contribution
of the topological penalty remains to be isolated, and it improves
sequential-learning retention on Llama-3-8B, where the log-ratio Omega-S
regulariser leaves more code capability after a sequential fine-tune than no
regularisation at all on 9 of 10 seeds (absolute HumanEval $0.173 \to 0.238$,
$+37.7\%$ relative, one-sided sign test $p = 0.011$), and beats tuned weight
decay on 10 of 10 seeds ($p = 0.002$) and tuned EWC on 8 of 10 ($p = 0.014$)
on the same absolute measure, without access to previous-task data.


Throughout, ``resilience'' names a structural configuration and not a
measured return time: no result below measures recovery from a perturbation,
and Remark~\ref{rem:nomenclature} states the scope of the term.

\textbf{Keywords:} topological entropy, non-equilibrium thermodynamics,
clustering coefficient, Prigogine minimum entropy production, scale-free
networks, structural monopolies, neural network regularisation.
\end{abstract}

\tableofcontents
\vspace{1em}

% ============================================================
\section{Introduction}
% ============================================================

\subsection{The Structural Monopoly Problem}

Complex adaptive systems, among them biological ecosystems, financial
markets, power grids and neural networks, share a common failure mode:
\emph{structural monopoly}. A small subset of nodes or weights
accumulates disproportionate influence over the network's information
flow, creating fragility that absolute-count metrics fail to detect.
Diversity metrics count species, neurons, or assets without measuring
their functional interdependence. A network can appear diverse while
being structurally dominated by a single hub. That gap is not hypothetical:
a survey of 1\,921 soil samples across 33 countries reports that agricultural
inputs raise apparent alpha diversity while shifting communities toward less
mutualistic and more fragmented configurations, and states explicitly that
higher diversity need not imply better function when the added taxa are
opportunistic rather than functionally redundant~\cite{saati2026}.

This paper addresses the following question: \emph{what is the correct
mathematical quantity to measure structural monopoly, and what is its
thermodynamic meaning?}

\subsection{The Central Object: Tr$(A^3)$}

The clustering coefficient of a graph with adjacency matrix $A$ is
standardly defined as~\cite{watts1998}:
\begin{equation}
C = \frac{\operatorname{Tr}(A^3)}{\sum_{ij}(A^2)_{ij} - \operatorname{Tr}(A^2)}.
\label{eq:clustering}
\end{equation}

The numerator $\operatorname{Tr}(A^3) = \sum_{ijk} A_{ij}A_{jk}A_{ki}$
counts the number of closed triangles (3-cycles) in the graph. For
weighted matrices $W \in \mathbb{R}^{m \times n}$ arising from neural
network layers, we define the induced adjacency matrix
$A = WW^\top \in \mathbb{R}^{m \times m}$ and the topological pressure
functional:
\begin{equation}
\mathcal{T}(W) = \operatorname{Tr}(WW^\top)^3 = \operatorname{Tr}(A^3).
\label{eq:T_functional}
\end{equation}

We will show that minimising $\mathcal{T}(W)$ is equivalent to
minimising structural monopoly, and that this minimisation corresponds
to a Prigoginian minimum-entropy-production principle.

\subsection{Relation to Prior Work}

The connection between triangle density and phase transitions in
networks was established by Strauss~\cite{strauss1986} via the
exponential random graph model with triangle fugacity $\lambda$.
Recent work~\cite{motif2026} shows that above a critical fugacity,
random graphs undergo a first-order transition from a triangle-poor
to a triangle-rich phase, accompanied by a change in spectral density
of the adjacency matrix, the same object we study here.

The connection between network topology and non-equilibrium
thermodynamics has been partially explored by Topologically-constrained
fluctuations~\cite{martins2023} and thermodynamic resilience
theory~\cite{toporesilience2017}. Our contribution is to show that
$\mathcal{T}(W) = \operatorname{Tr}(A^3)$ is the \emph{correct}
topological entropy production rate whose minimisation implements
Prigogine's theorem~\cite{prigogine1967}.

AlphaPruning~\cite{lu2024alpha} uses heavy-tailed spectral exponents
of weight matrices to guide neural network pruning. It is natural to read
that approach as ``spectral'' in contrast to a ``topological'' one based on
triangle counting, but Proposition~\ref{prop:spectral_identity} shows the
contrast is not available at the level of the global functional:
$\operatorname{Tr}(A^3)$ \emph{is} the third spectral moment. The two
approaches read different functionals of the same spectrum. Where our
framework does depart is (i) in supplying a thermodynamic interpretation of
why such structural signals predict redundancy, (ii) in being a
\emph{training-time gradient flow} rather than a post-hoc statistic, and
(iii) in the node-wise field $\operatorname{diag}(A^3)$, which by
Proposition~\ref{prop:node_nonspectral} is not spectrally determined.

% ============================================================
\section{Mathematical Framework}
% ============================================================

\subsection{The Functional Symbiotic Resilience Index}

\begin{definition}[FSRI]
\label{def:fsri}
For a network with adjacency matrix $A \in \mathbb{R}^{n \times n}$,
the Functional Symbiotic Resilience Index is:
\begin{equation}
\Omega = \frac{C \cdot D}{M \cdot \mathrm{Coex}},
\label{eq:fsri}
\end{equation}
where $C$ is the global clustering coefficient~\eqref{eq:clustering},
$D = \frac{\sum_{ij} A_{ij}}{n(n-1)}$ is the connection density,
$M = 1/\lambda_2(L)$ is the inverse algebraic connectivity, where
$\lambda_2$ is the second-smallest eigenvalue of the graph Laplacian
$L = \operatorname{diag}(\mathbf{k}) - A$, and
$\mathrm{Coex} = \mathrm{Var}(\mathbf{k})$ is the variance of the degree
sequence $\mathbf{k} = A\mathbf{1}$.

Earlier statements of this definition left $M$ as ``a modularity measure, for
instance an inverse spectral gap''. We fix it here because that latitude turned
out to matter: $M$ sits in the denominator, so the choice between $\lambda_2$
and $1/\lambda_2$ \emph{inverts} how $\Omega$ responds to compartmentalisation.
The orientation is not a matter of convention. It is fixed by the observation
the index was built to formalise, which Remark~\ref{rem:factors} sets out: a
graph that separates easily into modules has small $\lambda_2$, so $1/\lambda_2$
is the quantity that is large on modular graphs, and it is modularity, not
connectivity, that must appear in the denominator.
\end{definition}

\begin{figure}[ht]
\centering
\includegraphics[width=0.80\textwidth]{figure1_monopoly.pdf}
\caption{What the index responds to, at matched edge count. Left: connectivity
distributed across the network. Right: a structural monopoly, in which a hub
and a tight core absorb most of the closed three-step walks. Node labels give
the triangles incident on each node, $(A^3)_{ii}/2$, which is the quantity the
node-wise decomposition of Section~\ref{sec:spectral} reads off.}
\label{fig:monopoly}
\end{figure}

\begin{remark}
\label{rem:factors}
$\mathrm{Coex}$ (Competitive Exclusion) measures degree heterogeneity.
In the ecological interpretation, which derives from prior work on
co-inclusion and co-exclusion networks in agricultural
microbiomes~\cite{biome_patent,acedo2022indices}, high $\mathrm{Coex}$
indicates that a few nodes dominate connectivity, analogous to competitive
exclusion in resource competition. The FSRI $\Omega$ is large when the
network is simultaneously clustered ($C$ high), dense ($D$ high),
cohesive rather than compartmentalised ($M = 1/\lambda_2$
low, so the graph does not separate easily into modules), and egalitarian
($\mathrm{Coex}$ low).

That combination is not chosen for algebraic convenience: it is the one
reported for the systems the index was built to identify. In a survey of 350
vineyard soils across two countries~\cite{ortiz2021}, communities under
low-intervention management, which are the ones associated with small-world
structure and with resistance to perturbation, show simultaneously
\emph{higher} clustering, \emph{lower} modularity and \emph{lower} co-exclusion
than conventionally managed soils, with intermediate management intermediate on
all three. Every factor of~\eqref{eq:fsri} inherits its direction from that
observation, and the orientation of $M$ inherits it too: for $\Omega$ to be
large on those communities it must decrease with modularity, and since $M$ is
in the denominator, $M$ must be a modularity measure. An earlier statement of
this remark used the word ``modular'' where it should have read
``compartmentalised'' and pointed the other way; that was an error of phrasing,
not of the formula.

A second published study bears on the same direction at a larger scale. In a
global survey of soil microbiomes~\cite{saati2026}, agricultural inputs raise
modularity and shift network organisation toward compartmentalised
configurations that the authors read as less mutualistic and more fragmented.
We rest on the modularity and clustering results of that work and not on its
density result, because the density contrast there is drawn between groups of
very unequal size and the number of samples enters the network construction
twice, through a prevalence filter and through the power of the association
test; a comparison of density between groups of unequal sample size therefore
carries a component of sampling effort, in the sense made precise in
Section~\ref{sec:limitations}.

One point of nomenclature should be made explicit, since the name is
inherited but the quantity is not. In the microbiome work referred to
above, the co-exclusion quantity is the \emph{proportion of negative
(co-exclusion) edges} in a signed association network, measured
alongside clustering and modularity computed on the positive
subnetwork. Here $\mathrm{Coex}$ is the variance of the degree sequence
of a single unsigned network. These are different measurements, and the
ecological reading transfers as an analogy, not as the same statistic.
\end{remark}

\begin{remark}[Nomenclature: what ``resilience'' names here]
\label{rem:nomenclature}
The name of the index is inherited from the applied setting in which it was
first used, and it is worth separating from the technical sense the word
carries in ecology. Engineering resilience is a return time to equilibrium
after a perturbation, and ecological resilience in Holling's sense is the
magnitude of perturbation a system absorbs before changing basin. \emph{No
quantity in this paper measures either.} What $\Omega$ measures is a
structural configuration: clustered, dense, weakly compartmentalised and
egalitarian in degree. The link from that configuration to resilience in
either technical sense is an inference from the literature of
Remark~\ref{rem:factors}, not a measurement reported here, and where the
inference has been tested it has been tested on function rather than on
recovery~\cite{wagg2019}. We keep the name because it is the one under which
the index is used and cited, and we state the gap rather than let the name
assert it. The same caution applies to ``health'' as a synonym: soil health
is a policy-defined term tied to one domain, and adopting it would narrow an
index that Section~\ref{sec:applications} applies to weights, prices and
grids.
\end{remark}

\begin{remark}[Scope condition on the orientation of $M$]
\label{rem:modularity_scope}
Of the four factors, $M$ is the one whose \emph{orientation} is least settled
by evidence, and we state the condition rather than leave the choice of
Definition~\ref{def:fsri} looking universal. Three independent lines point the
other way from Remark~\ref{rem:factors}. Long-term experimental warming of
grassland soil raises network size, connectivity, clustering \emph{and}
relative modularity together, and the resulting networks are measured to be
more robust, not less~\cite{yuan2021}. In trophic settings a modular and
weakly connected architecture is the stabilising one, while dense and nested
architecture stabilises mutualistic networks, so the sign of modularity is a
function of which interaction type dominates rather than a constant of
nature~\cite{thebault2010}. And in the human gut, co-occurrence networks of
healthy subjects are reported to be \emph{more} easily disrupted by removal of
a few key taxa than those of patients with inflammatory bowel
disease~\cite{baldassano2016}.

None of these contradicts the direction adopted here, because none is measured
on the system Remark~\ref{rem:factors} describes: a competitive and
mutualistic microbial community under a management-intensity gradient. What
they establish is that the orientation of $M$ is domain-conditional. We
therefore restrict Definition~\ref{def:fsri} to settings in which competitive
and mutualistic interactions dominate and compartmentalisation reflects niche
specialisation rather than the buffering of a trophic cascade, and we note
that in the two applied settings of this framework the question does not
arise: under the construction of Remark~\ref{rem:composition} the modularity
factor is numerically inert, so the applied results are indifferent to its
orientation. Modularity is not thereby demoted: in the domain of origin it is
among the strongest discriminators available. It is the sign, not the
information content, that is conditional.
\end{remark}

\begin{remark}[Relation to the applied implementation]
\label{rem:composition}
Definition~\ref{def:fsri} fixes the index. The regulariser of the companion
applied work~\cite{acedo2026omegas} computes a related but distinct objective:
it differs in the construction of $A$, in the normalisation of $C$ and in the
orientation of $M$, and that paper documents each difference and measures which
of the four factors is live under its construction. We record the relation here
rather than the detail, because no result of the present paper depends on it:
Propositions~\ref{prop:spectral_identity},~\ref{prop:node_nonspectral},
\ref{prop:gradient_flow} and~\ref{prop:shape_invariant}, and the curvature
bound of Section~\ref{sec:orc_graphs}, are statements about
$\operatorname{Tr}(A^3)$ and about individual factors, not about how the four
are composed. A financial companion paper places $\mathrm{Coex}$ on the other
side of the ratio, and that difference is now settled rather than pending, and
is a consequence rather than a convention. If degree heterogeneity is
destabilising in the sense of~\cite{yan2017}, then a high $\mathrm{Coex}$ marks
fragility in both domains, and the two placements are two packagings of one
claim: an index built to be large on resilient systems carries the
concentration factor in the denominator, and an indicator built to rise under
stress carries it in the numerator. Under financial stress every factor moves the same
way, correlations rising so that $C$, $D$ and $\mathrm{Coex}$ all increase while
the inverse spectral gap falls, so an indicator meant to rise under stress must
compose them differently from an index meant to be large on resilient systems.
The asymmetry beneath that is worth stating here because it bears on the scope
of Definition~\ref{def:fsri}: a densely interconnected, weakly modular community
is the healthy configuration in the ecological setting of
Remark~\ref{rem:factors}, and is the crisis configuration in a market. The index
as defined is oriented for the former, and applying it unchanged to the latter
would call a crash resilient.

One point does belong here, because it bears on the scope of the index rather
than on any implementation. Under the bounded map that the applied work uses to
build $A$ from weight matrices, three of the four factors become numerically
inert: their elasticity with respect to the weights is at or below $10^{-4}$,
against $9\times10^{-3}$ for $\mathrm{Coex}$. That is a property of \emph{that
construction} and not of the index. In the ecological setting from which the
index derives, modularity is among the strongest discriminators available: in
the survey of Remark~\ref{rem:factors} it separates management types with
$r^2 = 0.955$, comparable to co-exclusion and higher than clustering. A factor
that carries most of the signal in the domain of origin and none under a
saturating construction on neural weights is telling us about the construction,
which is the reading Remark~\ref{rem:doubly_stochastic} already develops for a
different degenerate case.
\end{remark}

\begin{remark}[Conditions of comparability across graphs]
\label{rem:comparability}
The four factors do not behave alike under a change in the order of the
graph, and this constrains when two values of $\Omega$ may be compared.
$C$ is a ratio of closed triangles to connected triples, and is
dimensionless and bounded in $[0,1]$. $M$ depends on which modularity
measure is chosen, and neither standard choice is size-invariant a priori:
a partition-quality score is bounded but inherits the resolution limit
of~\cite{fortunato2007}, under which the score of a graph with fixed
community structure drifts with $n$, while an inverse spectral gap is
unbounded and its scaling with $n$ depends on the graph family.
$D = \sum_{ij} A_{ij}/n(n-1)$ is $O(1/n)$ at fixed mean degree.
$\mathrm{Coex} = \mathrm{Var}(\mathbf{k})$ carries units of degree squared
and, for a degree distribution with tail exponent $\gamma < 3$, diverges
with $n$; throughout, $\gamma$ denotes that degree exponent, while the
sampling exponents of Section~\ref{sec:limitations} are written $\beta_V$
and $\beta_E$. The index therefore multiplies two size-invariant quantities by
two that are not, so a comparison of $\Omega$ between graphs of different
order measures order as well as structure, and
Proposition~\ref{prop:ws_exact} makes the dependence explicit for
Watts-Strogatz graphs.

A second and weaker condition should be kept separate from this one. Being
dimensionless and bounded does not make a factor invariant under a change
of \emph{sampling} effort, in the case where the graph is estimated rather
than given. Section~\ref{sec:limitations} reports a measurement in which
all four factors, $C$ included, drift under such a change. All comparisons
reported in this paper are at matched order and matched edge count
(Proposition~\ref{prop:discriminative} at $N = 50$,
Table~\ref{tab:pressure} at $N = 40$, and Figure~\ref{fig:monopoly} at
matched edge count), and we state the condition rather than leave it
implicit. A constructive counterpart, identifying which two of the
four factors carry the dependence and what removes it, is given in
Proposition~\ref{prop:shape_invariant}.

The condition can be made numerical rather than left as an injunction. If
$\Omega$ scales as $n^{b}$, then two graphs whose orders differ by a factor
$r$ differ by $r^{|b|}$ on that account alone, so admitting a tolerance
$\varepsilon$ of order-induced variation bounds the admissible ratio at
$r \le (1+\varepsilon)^{1/|b|}$. On the estimated graphs of
Section~\ref{sec:limitations}, expressed against the order rather than against
the sample size, $b = -2.18$ for $\Omega$ and $+0.05$ for
$\widetilde{\Omega}$. A ten per cent tolerance therefore admits orders
differing by about $4\%$ for $\Omega$ and by a factor of about $6$ for
$\widetilde{\Omega}$. The numbers are family-dependent and are quoted to fix
an order of magnitude, not as constants.

Finally, the condition is \emph{vacuous} wherever the node set is fixed by
construction rather than discovered by sampling. In the neural application of
Section~\ref{sec:omega_s} the nodes are the rows of a given weight matrix, and
in a correlation network over a fixed set of assets the nodes are the assets;
in neither case does $n$ vary with how much data is collected. Everything in
this remark and in Section~\ref{sec:limitations} concerns graphs whose vertex
set grows with the sampling effort, and the applied results of this framework
lie outside that class.
\end{remark}

\begin{proposition}[The clustering factor carries a density scale]
\label{prop:excess}
Let $A = c\,(J - I)$ with $c > 0$, that is a graph of constant affinity $c$
and zero diagonal, where $J$ is the all-ones matrix. Then
\begin{equation}
C = c, \qquad D = c, \qquad \text{hence} \quad C/D = 1
\label{eq:excess}
\end{equation}
identically, for every order $n \ge 3$ and every $c$.
\end{proposition}

\begin{proof}
With $A = c(J-I)$ one has $A^2 = c^2[(n-2)J + I]$ and
$\operatorname{Tr}(A^3) = c^3 n(n-1)(n-2)$, while
$\sum_{ij}(A^2)_{ij} - \operatorname{Tr}(A^2) = c^2 (n-2)n(n-1)$, so
$C = c$ by~\eqref{eq:clustering}. And $\sum_{ij} A_{ij} = c\,n(n-1)$ gives
$D = c$. \qed
\end{proof}

\begin{remark}[The density-free triadic excess, and how to calibrate it]
\label{rem:excess}
Proposition~\ref{prop:excess} is elementary but it has a practical
consequence that is easy to miss when $C$ is read on its own: a value of $C$
is not by itself evidence of triadic structure, because a graph with no
structure whatever returns $C$ equal to its mean affinity. The
density-free quantity is the ratio $C/D$, which equals unity on the
constant-affinity graph and departs from unity only in the presence of
triadic organisation. This gives the ratio an absolute reference point that
$C$ lacks.
The reference point can also be obtained empirically, without assuming
constant affinity, by permuting the entries of the matrix from which $A$ is
built: the permutation preserves the distribution of affinity values exactly
and destroys only their arrangement, so any excess of $C/D$ above the
permuted value is attributable to arrangement. In the weight-matrix
construction discussed below, twenty-eight such permutation nulls return
$C/D$ equal to $1.0000$ to four decimal places, in agreement
with~\eqref{eq:excess}.
\end{remark}

\begin{remark}[A construction under which the index is degenerate]
\label{rem:doubly_stochastic}
Two of the four factors are not merely size-dependent, as
Remark~\ref{rem:comparability} records, but can be made structurally
constant by a choice of construction, and it is worth stating the one case in
which this makes $\Omega$ inevaluable. Suppose $A$ is scaled to be doubly
stochastic with zero diagonal, as is standard practice in spectral
clustering. Then $\mathbf{k} = A\mathbf{1} = \mathbf{1}$, so
$\mathrm{Coex} = \mathrm{Var}(\mathbf{k}) = 0$ exactly; and
$\sum_{ij} A_{ij} = n$, so $D = 1/(n-1)$, a constant of the order alone.
The index~\eqref{eq:fsri} therefore has a vanishing denominator and is not
defined, while the well-defined content of that construction is the ratio
$C/M$. This is a scope condition on Section~\ref{sec:applications} rather
than a defect: it says that $\Omega$ presupposes a construction in which the
degree sequence is free to vary, and it identifies $C$ and $M$, which are
also the two factors that
Section~\ref{sec:limitations} finds least dependent on sampling effort, as
the pair that survives when it is not.
\end{remark}

\begin{remark}[Which factor moves when the index is used as a training penalty]
The four factors above are all well-defined properties of a given weighted
graph, and the index is descriptive as stated. A separate, narrower
observation concerns the \emph{specific} case where $\Omega$ is minimised as a
regulariser during LLM fine-tuning, under the pseudo-adjacency construction
$A=\sigma(|WW^\top|)$ used in the applied work~\cite{acedo2026omegas}. There, a
direct measurement finds that the clustering factor $C$ does not move, so a
gradient-based optimiser can only act on $\mathrm{Coex}$.

A subsequent measurement over twenty-eight attention projection matrices of
an eight-billion-parameter decoder model refines the reason, and the refined
statement is both stronger and more specific than the one we gave
previously. The off-diagonal entries of $|WW^\top|$ are small rather than
large, with a median near $0.07$, so the logistic map does not pin the
affinities near its ceiling; it collapses them onto the midpoint of its
range, and $C$ is measured at $0.51$ rather than at $1$. The consequence is
sharper than saturation. Read through
Proposition~\ref{prop:excess}, a clustering value equal to the mean affinity
is what a graph with no triadic organisation returns, and the measured excess
$C/D$ on that construction is $1.0000$, equal to its permutation null to four
decimal places, in every module tested. The factor is not pinned at a
ceiling: the constructed adjacency carries no triadic content over its
density at all. Consistently with this, the dimensionless sensitivity of
$\operatorname{Tr}(A^3)$ to $W$ under that construction is zero at working
precision in at least one of the modules measured.

Two further points complete this. The first is a discrepancy between how the
construction is described and how it is implemented, which the applied
work~\cite{acedo2026omegas} now reports and which we repeat here because the
description above inherits it: the reference implementation forms
$\sigma(|WW^\top|)$ only when $W$ is non-square, and applies $\sigma(|W|)$
elementwise when it is square. The measurements quoted above used one module
of each kind. The degree-variance factor differs by roughly a factor of $400$
between the two branches on the square module measured, so the two are not
interchangeable, and any reimplementation should fix the branch deliberately.

The second is that the contrast-preserving construction, which an earlier
version of this remark offered as the obvious repair, has since been built and
evaluated. A row-wise cosine affinity does what it was designed to do: the
measured excess $C/D$ rises from $1.0000$ to a median of $1.36$ across the same
twenty-eight modules, the clustering factor moves off the value predicted by
Proposition~\ref{prop:excess} for an unstructured graph, and the sign of the
penalty controls the direction of movement during training in every module
tested. And retention becomes \emph{worse}, on all ten seeds of the applied
protocol. The inert clustering channel was therefore not a defect suppressing
the method's effect, and we record the negative outcome rather than leave the
repair standing as an open suggestion. The consequence for this paper is
narrow but worth stating: the identification of $\mathcal{T}$ with structural
monopoly is a property of the functional, and the empirical fact that
activating the clustering channel does not improve the applied outcome bears
on the regulariser, not on Propositions~\ref{prop:spectral_identity}
and~\ref{prop:node_nonspectral}.
\end{remark}

\begin{remark}[Whether the operative factor reduces to row norms]
A second measurement in the applied work bears on the same point. Since the
operative factor there is $\mathrm{Coex}$, the variance of the degree sequence,
and a node's degree correlates strongly with the norm of the corresponding
weight row, the natural objection is that the effect reduces to equalising
those norms. A control that penalises only row-norm variance, swept on a grid
of its own rather than borrowed from another arm, does not reliably improve
retention over no regularisation at all, while the full index penalty exceeds
that control on 8 of 10 seeds. The factor carrying the effect is therefore
$\mathrm{Coex}$ as computed on the weight graph, and not the row norms with
which it is correlated, though at ten seeds that separation is suggestive
rather than conclusive.

Two further points should be recorded because they bound the objection more
precisely than the control arm does. First, the association is not incidental
to the construction but mechanical in it: under $A = \sigma(|WW^\top|)$ the
degree of a node contains the squared norm of its own row through the
diagonal of $WW^\top$, and the rank correlation between a node's degree and
its row norm, measured within modules on the same twenty-eight matrices, is
$+0.61$ in the median and positive in all of them, the smallest value being
$+0.10$. An association that holds in every module with a common sign is
structural, not a sampling artefact, and the honest reading is that on this
construction the concentration channel and the row-norm distribution are
close to interchangeable. Second, that mechanical component is a property of
the construction and not of the factor: an affinity construction in which
each node's own scale is divided out before the product is formed makes the
node's degree \emph{provably} invariant to the norm of its own row, since
rescaling a row leaves every pairwise affinity unchanged. Under such a
construction any residual association can only reflect covariation between
norm and direction in the weights themselves, and is measured to change sign.
We state the property rather than the construction, which belongs to the
applied work.
\end{remark}

\begin{remark}[Direction of the clustering factor]
\label{rem:direction}
A reader who substitutes~\eqref{eq:clustering} into~\eqref{eq:fsri} will
notice that $\Omega$ is linear in $\operatorname{Tr}(A^3)$ and carries it
in the \emph{numerator}, so lowering $\operatorname{Tr}(A^3)$ lowers
$\Omega$. Read literally, minimising the topological pressure would
therefore move a network toward the end of the index that the index
itself calls fragile. We state the tension rather than leave it for the
reader to find. It is resolved by observing that a high
$\operatorname{Tr}(A^3)$ is ambiguous: triangles distributed across the
graph give $C$ high and $\mathrm{Coex}$ low, hence $\Omega$ high,
whereas the same triangle count concentrated on a hub and its core gives
$C$ high and $\mathrm{Coex}$ high, hence $\Omega$ low. What
distinguishes the two configurations is $\mathrm{Coex}$, in the
denominator, and not $\operatorname{Tr}(A^3)$ on its own. In the applied
setting the composite objective is observed to reduce degree variance
substantially, so the equivalence between minimising
$\operatorname{Tr}(A^3)$ and suppressing structural monopoly rests on
that \emph{measured} reduction and not on the algebra of $\Omega$, which
by itself does not deliver it.

There is a second reading of the same tension, which we record because it
locates a limit of the composite rather than a defect of the functional. The
numerator of $\Omega$ treats triadic closure as a benefit and the penalty of
Section~\ref{sec:omega_s} treats it as a cost, and both readings are defensible
in their own regime. That is the structure of a two-sided constraint on
connectivity, and Remark~\ref{rem:window} sets out what the index as composed
does and does not express about it.
\end{remark}

\begin{proposition}[Discriminative Power]
\label{prop:discriminative}
On synthetic networks with $N = 50$ nodes and matched edge count, drawn from
the preferential-attachment~\cite{barabasi1999} and small-world~\cite{watts1998}
families and from the ring lattice, the FSRI
of Definition~\ref{def:fsri} orders topologies as
\[
\Omega_{\text{ring}} \to \infty \;\gg\; \Omega_{\text{WS}} \approx 0.049
\;\gg\; \Omega_{\text{ER}} \approx 0.0031 \;\gtrsim\;
\Omega_{\text{BA}} \approx 0.0018,
\]
a separation of one and a half orders of magnitude among the finite values.
\end{proposition}

The three finite entries are medians over forty seeds, with ranges
$[0.029, 0.092]$, $[0.0013, 0.0064]$ and $[0.0013, 0.0025]$ respectively; the
order between the two sparse random families is therefore not stable at the
level of an individual seed, while the separation of both from Watts-Strogatz
is. The separation is carried by the concentration factor, whose median values
across the same seeds are $0.58$, $4.88$ and $16.46$: the index is reading
degree concentration, which is what the proposition is meant to establish.

The ring lattice requires a word, because its entry is a divergence and not a
large number. A $K$-regular ring has $\mathrm{Coex} = \mathrm{Var}(\mathbf{k})
= 0$ identically, so the denominator of~\eqref{eq:fsri} vanishes and $\Omega$
is not defined. This is the degenerate case of Remark~\ref{rem:doubly_stochastic} and
it agrees with the divergence at $p \to 0$ derived independently in
Proposition~\ref{prop:ws_exact}. We record it as a divergence rather than
substituting a numerical floor for the vanishing variance, because the value
such a substitution returns is fixed by the floor and not by the graph. The
consequence for how the finite entries should be read is that they measure
distance from an unattainable regular ideal rather than position on a bounded
scale, and that the index saturates at the egalitarian end. That saturation is
the same property discussed in Remark~\ref{rem:window}, reached here from the
elementary direction.

\subsection{Topological Pressure and the Gradient of Tr$(A^3)$}

The key analytical result connecting $\mathcal{T}$ to structural
monopoly is:

\begin{theorem}[Topological Pressure Gradient]
\label{thm:pressure_gradient}
The gradient of the topological pressure functional with respect to the
adjacency matrix is:
\begin{equation}
\frac{\partial}{\partial A_{ij}} \operatorname{Tr}(A^3) = 3(A^2)_{ij}.
\label{eq:gradient}
\end{equation}
The entry $(A^2)_{ij} = \sum_k A_{ik}A_{kj}$ counts the number of
paths of length 2 between nodes $i$ and $j$, i.e., the number of
common neighbours. High $(A^2)_{ij}$ indicates that the edge $(i,j)$
participates in many triangles and therefore contributes
disproportionately to the clustering structure.
\end{theorem}

\begin{proof}
Direct computation: $\frac{\partial}{\partial A_{ij}}
\operatorname{Tr}(A^3) = \frac{\partial}{\partial A_{ij}}
\sum_{klm} A_{kl}A_{lm}A_{mk} = 3(A^2)_{ij}$,
using the symmetry of $A$ and the cyclicity of trace.
\end{proof}

\begin{definition}[Topological Pressure Distribution]
The \emph{topological pressure} on edge $(i,j)$ is $\pi_{ij} = 3(A^2)_{ij} \cdot A_{ij}$.
The \emph{topological pressure Gini coefficient} $G(\pi)$ measures
the inequality of this distribution across edges.
\end{definition}

Our experiments demonstrate (Table~\ref{tab:pressure}):

\begin{table}[ht]
\centering
\caption{Topological pressure distribution by network type ($N=40$).}
\label{tab:pressure}
\begin{tabular}{lcccc}
\toprule
Network type & Max pressure & Mean pressure & Gini ($G$) & $\mathcal{T}/N^3$ \\
\midrule
Barabási-Albert (monopoly) & 15.00 & 4.03 & 0.986 & 0.00478 \\
Ring lattice (uniform)     & 12.00 & 9.00 & 0.272 & 0.01125 \\
\midrule
Ratio (monopoly/uniform)   & 1.25× & 0.45× & \textbf{3.6×} & 0.42× \\
\bottomrule
\end{tabular}
\end{table}

The Gini coefficient of topological pressure is the sharpest
discriminator: monopoly networks have near-perfect pressure
concentration ($G = 0.986$), while uniform networks distribute
pressure evenly ($G = 0.272$).

\subsection{Hutchinson Estimation of $\mathcal{T}(W)$}

For large weight matrices, exact computation of $\operatorname{Tr}(A^3)$
requires $\mathcal{O}(N^3)$ operations. The Hutchinson stochastic trace
estimator~\cite{hutchinson1990} provides an unbiased $\mathcal{O}(N^2)$
approximation:
\begin{equation}
\widehat{\operatorname{Tr}}(A^3) = \frac{1}{n_p} \sum_{i=1}^{n_p}
\mathbf{z}_i^\top A^3 \mathbf{z}_i,
\label{eq:hutchinson}
\end{equation}
where $\mathbf{z}_i \sim \text{Rad}(\pm 1)$ are Rademacher probe vectors
and $A^3 \mathbf{z} = W(W^\top (W(W^\top (W(W^\top \mathbf{z})))))$
is computed without materialising $A$.

\subsection{Spectral Identity and the Locus of Non-Spectral Information}
\label{sec:spectral}

A question that determines how $\mathcal{T}$ should be positioned relative
to spectral methods is whether $\operatorname{Tr}(A^3)$ carries information
beyond the eigenvalue spectrum of $A$. The answer is negative for the global
functional and positive for its node-wise decomposition, and the distinction
is sharp.

\begin{proposition}[The global functional is spectral]
\label{prop:spectral_identity}
For symmetric $A$ with eigenvalues $\{\lambda_i\}$,
\begin{equation}
\operatorname{Tr}(A^3) = \sum_i \lambda_i^3 .
\label{eq:third_moment}
\end{equation}
\end{proposition}
\begin{proof}
Immediate from the spectral decomposition $A = V\Lambda V^\top$ with
$V$ orthogonal: $\operatorname{Tr}(A^3) = \operatorname{Tr}(V\Lambda^3V^\top)
= \operatorname{Tr}(\Lambda^3)$.
\end{proof}

$\mathcal{T}(W)$ is therefore \emph{exactly the third moment of the
spectrum}, a spectral statistic in the same family as the Absorption
Ratio (leading-eigenvalue mass), effective rank and the Vendi score
(spectral entropies). Any two of these are functionals of the same
$\{\lambda_i\}$ and cannot, in general, be expected to carry mutually
exclusive information. We state this explicitly because it corrects a
framing of ``topological'' as opposed to ``spectral'' that the
triangle-counting interpretation invites but that
Eq.~\eqref{eq:third_moment} forbids at the level of the global scalar.

The non-spectral content appears one level down. Writing
$A = \sum_k \lambda_k \mathbf{v}_k\mathbf{v}_k^\top$, the node-wise
decomposition of the topological pressure is
\begin{equation}
(A^3)_{ii} = \sum_k \lambda_k^3\, (v_k)_i^2 ,
\label{eq:node_decomp}
\end{equation}
which depends on the \emph{eigenvectors}, not on the spectrum alone.

\begin{proposition}[The node decomposition is not spectrally determined]
\label{prop:node_nonspectral}
There exist cospectral graphs $G_1 \not\cong G_2$ with
$\operatorname{Tr}(A_1^3) = \operatorname{Tr}(A_2^3)$ but
$\operatorname{diag}(A_1^3) \neq \operatorname{diag}(A_2^3)$ up to
permutation.
\end{proposition}
\begin{proof}
By explicit construction. Exhaustive search over graphs on six vertices
yields a cospectral pair with common spectrum
$\{-1.9032,\,-1,\,-1,\,0.1939,\,1,\,2.7093\}$ and, necessarily by
Proposition~\ref{prop:spectral_identity}, common
$\operatorname{Tr}(A^3) = 12$, yet node-wise triangle counts
$\operatorname{diag}(A_1^3) = (4,2,2,2,2,0)$ and
$\operatorname{diag}(A_2^3) = (2,2,4,4,0,0)$, which differ as multisets.
\end{proof}

\begin{figure}[ht]
\centering
\includegraphics[width=0.80\textwidth]{figure2_cospectral.pdf}
\caption{The counterexample of Proposition~\ref{prop:node_nonspectral}. Both
graphs have the same eigenvalues and therefore, by
Proposition~\ref{prop:spectral_identity}, the same
$\operatorname{Tr}(A^3) = 12$. Their node-wise decompositions differ as
multisets, so $\operatorname{diag}(A^3)$ is not a function of the spectrum.
The pair is located by exhaustive search in the accompanying code rather than
asserted, and node labels are computed rather than drawn.}
\label{fig:cospectral}
\end{figure}

\begin{corollary}[Locus of the contribution]
\label{cor:locus}
The global index $\mathcal{T}$ is spectrally redundant: it can be matched
by other spectral functionals and should not be expected to dominate them
as a scalar detector. The node-wise field $\operatorname{diag}(A^3)$ is
not a function of the spectrum, and is therefore the component of the
framework that no eigenvalue-based scalar can reproduce.
\end{corollary}

Corollary~\ref{cor:locus} makes a prediction that is testable in any domain
where the framework is applied as a detector: the global scalar should tie
sharpened spectral baselines rather than beat them, while the node-wise
attribution, being parameter-free unlike eigenvector attribution, which
requires committing to a number of components $k$, should do better when the
number of structural epicentres is unknown. We state it as a prediction here.
A companion paper applying the index to financial correlation networks reports
a pattern consistent with it; that work uses a different composition of the
same four factors from the one in Definition~\ref{def:fsri} and we are
reconciling the two before treating its numbers as evidence for this
corollary.

\subsection{Signed Extension: Structural Balance}

The construction above uses $A = |{\cdot}|$ throughout, which discards the
sign of the underlying interaction. On a signed adjacency matrix
$A \in [-1,1]^{n\times n}$ the same functional acquires a different
meaning:
\begin{equation}
\operatorname{Tr}(A^3) = 6\,(\#\text{balanced triangles} -
\#\text{unbalanced triangles}),
\end{equation}
since a 3-cycle contributes $+|A_{ij}A_{jk}A_{ki}|$ when the product of
signs is positive and $-|A_{ij}A_{jk}A_{ki}|$ otherwise.

\begin{figure}[ht]
\centering
\includegraphics[width=0.92\textwidth]{figure3_signed.pdf}
\caption{Structural balance. The left and centre triangles carry the same
edge magnitudes and differ only in one sign, yet
$\operatorname{Tr}(A^3)$ separates them completely. Taking $|A|$, as the
unsigned framework and every spectral index computed on $|A|$ do, collapses
both to the graph on the right and destroys the distinction.}
\label{fig:signed}
\end{figure} This is exactly
Heider's structural balance~\cite{cartwright1956}, and it is invisible to
any functional of $|A|$, including all spectral indices computed on
$|A|$ and the unsigned $\mathcal{T}$ used in the rest of this paper.
Two networks with identical $|A|$ and opposite balance structure are
indistinguishable to the unsigned framework. Signed $\mathcal{T}$ is thus
a strict generalisation carrying information no unsigned index can access;
its behaviour in domains where sign is meaningful (ecological co-occurrence
versus co-exclusion, hedging versus co-movement in finance) is left for
future work.

% ============================================================
\section{Non-Equilibrium Thermodynamic Interpretation}
% ============================================================

\subsection{Prigogine's Theorem and Topological Entropy Production}

Prigogine's theorem states that in a linear non-equilibrium system
with fixed external constraints, the stationary state minimises the
rate of entropy production~\cite{prigogine1967,klein1954}. A rigorous
proof via network Markov chains was given by~\cite{schnakenberg_nets},
who showed that in networks with fixed boundary conditions, steady
states satisfy a minimum entropy production principle through
Schnakenberg's cycle observables.

\begin{remark}[Scope of Prigogine's Theorem]
Prigogine's theorem holds strictly in the \emph{linear} regime of
irreversible thermodynamics (near equilibrium, Onsager relations
valid). Extensions to non-linear regimes are contested~\cite{prigogine_critical},
and the principle can fail for strongly driven systems. We apply it
here as a \emph{guiding analogy} in the linear regime of weight
updates near a training fixed point, not as a rigorous claim about
arbitrary training dynamics.
\end{remark}

\begin{definition}[Topological Entropy Production Rate]
\label{def:tep}
For a network with weight matrix $W$, the topological entropy production
rate is:
\begin{equation}
\dot{S}_{\mathrm{topo}} \propto \mathcal{T}(W) =
\operatorname{Tr}(WW^\top)^3.
\label{eq:entropy_prod}
\end{equation}
This quantity measures the rate at which topological structure is
sustained through triangle-closing events, the network-theoretic
analogue of dissipative structure formation in Prigogine's framework.
\end{definition}

The analogy is made precise as follows. In Prigogine's framework,
a system far from equilibrium self-organises into a structured state
(dissipative structure) that sustains itself through energy throughput.
In a network under preferential attachment (analogous to energy
abundance $E_\text{ext}$), hub nodes accumulate connections, creating
structured but fragile configurations. Prigogine's theorem in network
form~\cite{schnakenberg_nets} implies that the \emph{stationary}
(minimum perturbation) configuration minimises $\dot{S}_\text{topo}$.

\begin{proposition}[Gradient Flow as Prigoginian Dynamics]
\label{prop:gradient_flow}
The gradient flow:
\begin{equation}
\dot{W} = -\lambda_\Omega \nabla_W \mathcal{T}(W)
= -6\lambda_\Omega \, (WW^\top)^2 W
\label{eq:gradient_flow}
\end{equation}
drives the weight matrix toward configurations of lower topological
entropy production. In the linear regime near a stationary point,
this is equivalent to Prigogine's minimum entropy production principle
applied to the space of weight configurations.

Fixed points $\dot{W} = 0$ satisfy $(WW^\top)^2 W = 0$, which holds
when $WW^\top$ has zero spectral radius on the support of $W$,
i.e., weight matrices inducing triangle-free graphs (maximum entropy
for given connection density).
\end{proposition}

\begin{proof}
Write $A = WW^\top$. From Theorem~\ref{thm:pressure_gradient},
$d\operatorname{Tr}(A^3) = 3\operatorname{Tr}(A^2\,dA)$, and
$dA = dW\,W^\top + W\,dW^\top$, so
\[
d\operatorname{Tr}(A^3)
= 3\operatorname{Tr}(A^2\,dW\,W^\top) + 3\operatorname{Tr}(A^2 W\,dW^\top)
= 6\,\langle A^2 W,\, dW\rangle ,
\]
the two terms being equal by symmetry of $A$ and cyclicity of the
trace. Hence
\begin{equation}
\nabla_W \operatorname{Tr}(WW^\top)^3 = 6\,(WW^\top)^2 W .
\label{eq:grad_W}
\end{equation}
Note the order of the factors: for $W \in \mathbb{R}^{m \times n}$ the
product $(WW^\top)^2 W$ is $m \times n$ and therefore conformable with
$W$, whereas $W(WW^\top)^2$ is not defined unless $m = n$. The fixed
point condition follows from $(WW^\top)^2 W = 0$.
\end{proof}

\begin{remark}[Connection to Maximum Caliber]
Prigogine's minimum entropy production principle can be derived from
the Maximum Caliber variational principle~\cite{maxcaliber2018} for
near-equilibrium systems with two coupled flows. This provides a
statistical mechanical foundation: the gradient flow~\eqref{eq:gradient_flow}
is the most probable trajectory for a stochastic weight update process
constrained to minimise topological entropy production.
\end{remark}

\subsection{Phase Transition in Triangle Density}

Our experiments confirm the existence of a topological phase transition.
Starting from an Erdős-Rényi network ($N=60$, $p=0.08$) and
progressively adding triangles via the fugacity mechanism
of~\cite{motif2026}, we observe:
\begin{itemize}
\item $\mathcal{T}/N^3$ increases monotonically from $1.94 \times 10^{-4}$
      (zero triangles) to $1.48 \times 10^{-2}$ (200 triangles added),
      a $76\times$ increase.
\item The degree entropy $S = -\sum_k p_k \log p_k$ increases from
      $1.95$ to $2.70$ as triangles are added, reflecting growing
      degree heterogeneity.
\item This is consistent with the theoretical prediction that
      triangle-rich phases exhibit heavy-tailed degree distributions
      with power-law exponents~\cite{motif2026}.
\end{itemize}

\begin{remark}[On the Proportionality $\Omega \propto (S_\text{env}/E_\text{ext})^2$]
An earlier formulation of this framework by the author proposes
$\Omega \propto (S_\text{env}/E_\text{ext})^2$. Our experiments on Watts-Strogatz networks
yield $r = -0.407$, suggesting the relationship is not a simple proportionality.
We now derive the exact form analytically.
\end{remark}

\begin{proposition}[Exact FSRI Formula for Watts-Strogatz Networks]
\label{prop:ws_exact}
For a Watts-Strogatz network with $n$ nodes, mean degree $K$, and
rewiring probability $p \in (0,1)$, using the exact clustering formula
$C(K,p) = \frac{3K-6}{4K-2}(1-p)^3$~\cite{barrat2000} and the
variance approximation $\mathrm{Var}(k) \approx Kp(1-p)$ valid for
$n \to \infty$:
\begin{equation}
\Omega(K,p) = \frac{3K-6}{4K-2} \cdot \frac{(1-p)^2}{np},
\label{eq:omega_ws}
\end{equation}
where $D = K/(n-1)$, written as $K/n$ to leading order, and
$\mathrm{Coex} = Kp(1-p)$. This formula achieves
correlation $r = 0.9997$ against numerical simulation across
$p \in [0.05, 0.9]$, $K = 6$, $N = 100$ (Table~\ref{tab:ws_validation}).

\begin{table}[h]
\centering
\caption{Analytical against numerical $\Omega(K=6, N=100)$ in
Watts-Strogatz, each simulated value the mean over 200 graphs. Produced by
\texttt{colabs/fsri\_colab3\_orc\_ws.py} in the accompanying repository.}
\label{tab:ws_validation}
\begin{tabular}{cccc}
\toprule
$p$ & $\Omega$ analytical & $\Omega$ numerical & relative error \\
\midrule
0.05 & 0.098455 & 0.113740 & 13.4\% \\
0.10 & 0.044182 & 0.050033 & 11.7\% \\
0.20 & 0.017455 & 0.018315 & 4.7\% \\
0.30 & 0.008909 & 0.008928 & 0.2\% \\
0.50 & 0.002727 & 0.002925 & 6.8\% \\
0.70 & 0.000701 & 0.001375 & 49.0\% \\
0.90 & 0.000061 & 0.001088 & 94.4\% \\
\bottomrule
\end{tabular}
\end{table}

\begin{remark}[What the correlation does and does not establish]
\label{rem:ws_scope}
The quoted $r = 0.9997$ is computed on levels that span three orders of
magnitude, so it is dominated by the small-$p$ end of the range and is a
weak test of agreement elsewhere. On the logarithms of the same seven points
it is $0.9447$, and the last column of Table~\ref{tab:ws_validation} is the
honest summary: the closed form is accurate to within fifteen per cent for
$p \le 0.5$ and departs badly beyond it, reaching a factor of eighteen at
$p = 0.9$.

The departure is expected and is not a failure of~\eqref{eq:omega_ws}. It
comes from the variance approximation $\mathrm{Var}(k) \approx Kp(1-p)$,
which holds as $n \to \infty$; at $n = 100$ the finite-size correction grows
with $p$, and at $p = 0.9$ the rewired graph retains almost none of the
lattice structure the approximation assumes. Every use made of the formula
in this paper is in the lattice-like regime of small $p$, where the table
shows it to be accurate, and the divergence as $p \to 0$ that
Proposition~\ref{prop:ws_exact} relies on sits at the accurate end.
\end{remark}

Setting $E_\mathrm{ext} = K$ (mean degree) and
$S_\mathrm{env} = \sqrt{Kp(1-p)}$ (degree standard deviation), so that
$(S_\mathrm{env}/E_\mathrm{ext})^2 = p(1-p)/K$ and hence
$p(1-p) = S_\mathrm{env}^2/E_\mathrm{ext}$, equation~\eqref{eq:omega_ws} can be
written exactly in terms of the thermodynamic variables. Multiplying
numerator and denominator by $(1-p)$,
\begin{equation}
\boxed{\;\Omega \;=\; \frac{3K-6}{4K-2} \cdot \frac{(1-p)^3}{n}
\cdot \frac{E_\mathrm{ext}}{S_\mathrm{env}^2}\;}
\label{eq:exact_relation}
\end{equation}
which reproduces~\eqref{eq:omega_ws} identically.

This shows three key results:
\begin{enumerate}
\item The ratio enters \emph{inverted}: $\Omega \propto E_\mathrm{ext} /
  S_\mathrm{env}^2$ and not $S_\mathrm{env}^2/E_\mathrm{ext}^2$ as the earlier
  formulation proposes. At fixed connectivity, more environmental stress
  \emph{lowers} the index rather than raising it, and higher connectivity at
  fixed stress raises it. The earlier formulation has the dependence the right
  way round in neither variable.
\item $\Omega$ still diverges as $p \to 0$ (pure lattice, maximum clustering),
  consistent with the ring lattice having $\Omega \to \infty$ in our
  experiments: the divergence is now carried by $S_\mathrm{env}^2 = Kp(1-p)$
  in the denominator, which vanishes with $p$.
\item The sign of the discrepancy with the earlier formulation is what the
  measurements show. Because the ratio belongs in the denominator, the
  empirical correlation between $\Omega$ and $(S_\mathrm{env}/E_\mathrm{ext})^2$
  is \emph{negative} ($r = -0.407$ to $-0.573$) rather than weakly positive.
  A proportionality would not produce a negative correlation of that size;
  an inverse relation does.
\item $\Omega \propto 1/n$ at fixed $K$ and $p$. The factor is visible
  in~\eqref{eq:omega_ws} and is not an artefact of the Watts-Strogatz
  parameterisation: it enters through $D = K/n$, which falls with the order
  of the graph at fixed mean degree, and it is reinforced by
  $\mathrm{Coex}$ whenever the degree distribution is heavy-tailed. The
  formula therefore delimits the index as a \emph{comparative} statistic in
  the sense of Remark~\ref{rem:comparability}: values of $\Omega$ obtained
  on graphs of different order are not directly comparable, and every
  comparison reported in this paper is at matched $n$.
\end{enumerate}

The factor $1/n$ isolated above is not irreducible: it is carried by two of
the four factors, and both have a counterpart that does not carry it.

\begin{definition}[Shape-normalised index]
\label{def:shape}
Let $G_k$ denote the Gini coefficient of the degree sequence
$\mathbf{k} = A\mathbf{1}$ and $\bar{k}$ its mean. The shape-normalised
index is
\begin{equation}
\widetilde{\Omega} = \frac{C \cdot \bar{k}}{M \cdot G_k},
\label{eq:omega_tilde}
\end{equation}
obtained from~\eqref{eq:fsri} by replacing the connection density with the
mean degree and the degree variance with the Gini coefficient of the same
degree sequence.
\end{definition}

\begin{proposition}[The shape-normalised index is order-invariant on
Watts-Strogatz graphs, at fixed $M$]
\label{prop:shape_invariant}
For Watts-Strogatz graphs with $n$ nodes, mean degree $K$ and rewiring
probability $p$, and with the modularity factor $M$ held fixed, the quantity
$\widetilde{\Omega}$ of Definition~\ref{def:shape} is a function of $K$ and
$p$ alone and does not depend on $n$, whereas $\Omega$ carries an explicit
factor $1/n$ by Proposition~\ref{prop:ws_exact}. The hypothesis on $M$ is not
cosmetic and Remark~\ref{rem:residual_drift} measures what it costs.
\end{proposition}

\begin{proof}
Each factor of~\eqref{eq:omega_tilde} is independent of $n$. The clustering
coefficient is $C(K,p) = \frac{3K-6}{4K-2}(1-p)^3$~\cite{barrat2000},
a function of $K$ and $p$. The mean degree is $\bar{k} = K$ by construction.
The rewiring process acts independently on each edge with probability $p$,
so the degree distribution of the Watts-Strogatz ensemble is determined by
$K$ and $p$ and not by the order of the graph; $G_k$, being a functional of
that distribution, is therefore also a function of $K$ and $p$. Taking $M$
constant, as in Proposition~\ref{prop:ws_exact}, gives
$\widetilde{\Omega} = \widetilde{\Omega}(K,p)$. By contrast
$D = K/(n-1)$ carries $1/n$ explicitly, which is the origin of the factor
$1/(np)$ in~\eqref{eq:omega_ws}.
\end{proof}

\begin{remark}[Why these two substitutions and not others]
The proof uses two properties that hold beyond the Watts-Strogatz family.
The Gini coefficient is scale-invariant, $G_k(c\mathbf{k}) = G_k(\mathbf{k})$
for $c>0$, whereas $\mathrm{Var}(c\mathbf{k}) = c^2\mathrm{Var}(\mathbf{k})$
and carries units of degree squared. The mean degree is intensive, whereas
$D = \bar{k}/(n-1)$ is not. The substitution therefore replaces two
quantities that depend on the scale of the degree sequence and on the order
of the graph by two that depend only on the shape of that sequence. This is
the general reason; Proposition~\ref{prop:shape_invariant} is the case in
which it can be checked exactly. Numerically, on Watts-Strogatz graphs with
$K=6$, $p=0.1$ and $n$ from $200$ to $6400$, the fitted slope of
$\log_{10}\Omega$ against $\log_{10} n$ is $-0.987$, consistent with
Proposition~\ref{prop:ws_exact}, while that of $\widetilde{\Omega}$ is
$+0.009$. Both figures hold $M$ fixed, as
Proposition~\ref{prop:shape_invariant} does.
\end{remark}

\begin{remark}[The residual drift, and where it comes from]
\label{rem:residual_drift}
Holding $M$ fixed is an idealisation, and on this family it is a false one:
the algebraic connectivity of a Watts-Strogatz graph falls with its order, so
$M = 1/\lambda_2(L)$ grows with $n$ and neither index is exactly
order-invariant. Measured over forty graphs per point at $K=6$, $p=0.1$ and
$n$ from $200$ to $1600$, with $M$ evaluated rather than fixed, the fitted
slopes become $-1.160$ for $\log_{10}\Omega$ and $-0.160$ for
$\log_{10}\widetilde{\Omega}$, against $-1.009$ and $-0.010$ for the same
graphs with $M$ held fixed. The whole of the difference is $\lambda_2$, whose
own slope is $-0.151$ on the same graphs; $C$, $\bar{k}$ and $G_k$ are flat to
within a few parts in a hundred across the range.

Two things follow. The substitution does most of what it is introduced to do,
removing about seven eighths of the drift: a factor of roughly $14$ per decade
of $n$ becomes a factor of roughly $1.4$. And it does not remove all of it,
because the concentration factor is not the only one that carries the order of
the graph. Restoring exact comparability would require a modularity factor
that is itself order-invariant, which we do not have and do not claim; the
comparisons reported in this paper are at matched $n$ for that reason.
\end{remark}

\begin{remark}[A third candidate for the concentration factor, with a
dynamical rather than a dimensional justification]
\label{rem:xi}
The argument above justifies $G_k$ on grounds of scale: it is invariant where
$\mathrm{Var}(\mathbf{k})$ is not. That argument says why the substitution
restores comparability; it does not say why a concentration measure belongs in
the denominator at all, which Remark~\ref{rem:factors} supplies only by
ecological analogy. A third candidate carries both. Write
\begin{equation}
\xi \;=\; \frac{\langle k^2 \rangle}{\langle k \rangle^2}
\;=\; 1 + \frac{\mathrm{Var}(\mathbf{k})}{\bar{k}^2}
\;=\; 1 + \frac{\mathrm{Coex}}{\bar{k}^2},
\label{eq:xi}
\end{equation}
the degree heterogeneity of the graph. It is $\mathrm{Coex}$ normalised by the
square of the mean degree, hence dimensionless and scale-invariant like $G_k$;
and it is the quantity for which degree heterogeneity is shown to destabilise
ecological networks with competitive and mutualistic interactions, entering a
stability criterion of the form
\begin{equation}
\sigma \sqrt{(S-1)\,\mathcal{C}\,\xi} \;<\; d
\label{eq:yan}
\end{equation}
for a community matrix with $S$ species, self-regulation $d$ and interaction
strengths of dispersion $\sigma$~\cite{yan2017}. We therefore record
$\widetilde{\Omega}_\xi = C\bar{k}/(M\xi)$ as a variant of
Definition~\ref{def:shape} on the same footing as the one stated there, and
note that the two agree in direction while differing in how they weight the
tail of the degree distribution.

Two cautions attach to importing~\eqref{eq:yan}, and both are about scope
rather than about $\xi$. The first is notational and easy to misread: in the
stability literature $\mathcal{C}$ denotes \emph{connectance}, the fraction of
realised links, which is the quantity written $D$ here, whereas $C$ in this
paper is the clustering coefficient. Translated into the notation of
Definition~\ref{def:fsri}, the criterion reads
$\sigma\sqrt{(n-1)\,D\,\xi} < d$; it says nothing about triadic structure. The
second is substantive. In that criterion connectance and heterogeneity enter
as a \emph{product}, and increasing either is destabilising, whereas
Definition~\ref{def:fsri} places $D$ in the numerator and the concentration
factor in the denominator, so the two are composed as a ratio and move
$\Omega$ in opposite directions. This is the long-standing tension between the
analytical stability criteria descending from~\cite{may1972,allesina2012} and
the empirical finding that connectivity accompanies function and
robustness~\cite{wagg2019,morrien2017,yuan2021}, and the two are not
measuring the same notion of stability: the criterion concerns local
asymptotic stability of a randomly parameterised community matrix, whereas the
evidence behind the numerator of~\eqref{eq:fsri} concerns functional
redundancy measured on ecosystem processes. We take from~\cite{yan2017} the
justification of the concentration factor and its normalisation, and not the
composition of the index, and we state the difference rather than let the
citation imply an endorsement it does not carry.
\end{remark}

\begin{remark}[Connectivity as a two-sided constraint, and what the index does
not express]
\label{rem:window}
Three independent lines converge on a statement the index as composed does not
make. In the stability criteria discussed above, connectance enters as a
destabilising quantity~\cite{may1972,allesina2012,yan2017}. In communication
graphs among language-model agents, reliability imposes a minimum degree and
error propagation a maximum, so that the admissible graphs occupy an interval
that may be empty~\cite{niu2026}. And in the ecological literature from which
this framework derives, the same shape has a name: Ulanowicz's window of
vitality, in which a flow network persists only in an intermediate band between
too little organisation, which leaves it incoherent, and too much, which leaves
it brittle for want of redundant pathways~\cite{ulanowicz2009}.

Definition~\ref{def:fsri} is monotone in each of its factors, and in $D$ in
particular. It therefore expresses one edge of such a window and not the window
itself: it says which direction is the resilient one, not where the direction
reverses. This is a limitation of the composite and not of the underlying
functional, since nothing in Propositions~\ref{prop:spectral_identity}
or~\ref{prop:node_nonspectral} depends on how the four factors are combined.
It also identifies the natural extension: an
index with an interior optimum in connectivity would require a second term
penalising excess, of the kind that the overhead side of the ascendency
decomposition supplies. Appendix~\ref{app:ascendency} carries that transplant
out exactly for the setting of this paper, and reports both what it delivers, a
closed form and a point prediction for the optimal degree, and what it does
not, since the resulting order parameter turns out to be a functional of the
degree sequence alone and therefore blind to wiring. We record the convergence because the three sources
share no method, and because agreement across epidemiology, random-matrix
stability and information-theoretic ecology is the kind of evidence a purely
structural argument cannot supply on its own.

One consistency check belongs here, because it delimits the reading rather than
weakening it. The four families of Proposition~\ref{prop:discriminative} are
compared at matched order and matched edge count, hence at matched mean degree,
and an information-theoretic order parameter computed on them separates them
hardly at all. The comparison that Proposition~\ref{prop:discriminative} reports
is therefore one of arrangement at fixed connectivity, which is the regime in
which the index is well posed. The limitation stated above concerns comparisons
made \emph{across} levels of connectivity, and no result of this paper makes
one.
\end{remark}

\begin{remark}[When the shape-normalised index is exact, and how to check it
in advance]
\label{rem:densification}
Proposition~\ref{prop:shape_invariant} replaces $D$ by $\bar{k}$ on the
grounds that the mean degree is intensive. That premise is not automatic on a
growing graph. If the edge count scales with the order as $m \sim n^{\delta}$,
then $\bar{k} = 2m/n \sim n^{\delta - 1}$ and $D \sim n^{\delta-2}$, so the
substitution is exact precisely when $\delta = 1$ and leaves a residual
exponent of order $\delta - 1$ otherwise. Real growing graphs are observed to
densify, with $\delta$ typically between $1$ and $2$~\cite{leskovec2005}, so
the residual is generally not zero: on the syntactic networks of
Section~\ref{sec:limitations} the two accumulation exponents give
$\delta = \beta_E/\beta_V = 0.955/0.77 = 1.24$, hence a predicted residual of
$+0.24$ per decade of order in $\bar{k}$, which is what is measured there.

The practical consequence is that $\widetilde{\Omega}$ carries its own
applicability condition, and one that can be checked before the index is
applied and without any additional data: $\delta$ is the slope of
$\log m$ against $\log n$ across the units to be compared, and both quantities
are already needed to compute $D$. Where $\delta$ is close to one the
substitution removes the order dependence almost entirely; where the graph
densifies strongly the substitution over-corrects, replacing a factor that
falls with $n$ by one that rises with it, and the shape argument of
Remark~\ref{rem:shape_scope} should not be expected to hold. We state the
condition as a prediction of the scaling argument; a systematic measurement
across families with different $\delta$ is left open.
\end{remark}

\begin{remark}[What Proposition~\ref{prop:shape_invariant} does not cover]
\label{rem:shape_scope}
Watts-Strogatz graphs have homogeneous degrees. The proposition establishes
invariance to the order of a given graph; it does not establish invariance
under a change of \emph{sampling} effort when the graph is estimated and its
degree distribution is heavy-tailed, which is the situation measured in
Section~\ref{sec:limitations}. The analytical and the empirical statements
agree in direction and are not the same claim. Nor does
$\widetilde{\Omega}$ inherit the ecological reading of $\mathrm{Coex}$ as a
degree variance: the Gini coefficient is a concentration measure of the same
sequence, so the reading transfers as one of inequality rather than of
variance.
\end{remark}

\end{proposition}

% ============================================================
\section{Curvature Interpretation: Triangles and Ollivier-Ricci Bounds}
\label{sec:orc_graphs}
% ============================================================


Before formalising that connection it is worth naming the closest
application of this curvature to the problem the framework addresses. Sandhu,
Georgiou and Tannenbaum~\cite{sandhu2016} compute Ollivier-Ricci curvature on
financial correlation networks as an indicator of market fragility and systemic
risk: the same construction and the same purpose as the financial embodiment of
this family, reached from discrete geometry rather than from triadic structure.
A subsequent analysis~\cite{expost2024} argues that the indicator identifies the
size and duration of a crisis accurately but does not lead it, that is that it
is an ex-post descriptor rather than an early warning. The same curvature has
since been applied to market instability with an explicitly network-geometric
reading~\cite{samal2021}, and, outside finance, as the basis of a community
detection method~\cite{sia2019}, which is a second use of the same local
quantity for a global structural purpose. Both bear on what
follows. The bound of Theorem~\ref{thm:orc_upper} relates the same
common-neighbour counts that the curvature reads, so the two are not
independent measurements of a network; and the ex-post limitation is a property
of that reading which a triadic reading does not automatically escape. We note
also that the relation between Ollivier-Ricci curvature and the spectrum of the
normalised Laplacian is established~\cite{bauer2012}, which places the curvature
of this section in the same spectral neighbourhood as
Proposition~\ref{prop:spectral_identity} rather than outside it.

The connection to differential geometry is formalised through the
\emph{Ollivier-Ricci curvature}~\cite{ollivier2009}, a discrete
analogue of Riemannian sectional curvature for metric spaces:
\begin{equation}
\kappa(x, y) = 1 - \frac{W_1(\mu_x, \mu_y)}{d(x,y)},
\label{eq:ricci}
\end{equation}
where $W_1$ is the Wasserstein-1 (earth-mover) distance between
the probability measures $\mu_x$, $\mu_y$ over neighbours of $x$
and $y$, and $d(x,y)$ is the graph distance.

\begin{proposition}[ORC and Triangle Density]
\label{prop:orc_triangles}
For a graph $G$ with adjacency matrix $A$ and uniform random-walk
measures $\mu_x(y) = A_{xy}/\deg(x)$, the Ollivier-Ricci curvature of
an edge $(x,y)$ is controlled by the number of common neighbours
$t_{xy} = |N(x) \cap N(y)|$, which is exactly the contribution of that
edge to $\operatorname{Tr}(A^3) = \sum_{(x,y) \in E} 2\,t_{xy}$. The
corresponding upper bound is Theorem~\ref{thm:orc_upper}; the condition
under which the curvature is non-negative is
Proposition~\ref{prop:orc_nonneg}.
\end{proposition}

\begin{theorem}[ORC Upper Bound]
\label{thm:orc_upper}
Let $t_{xy} = |N(x) \cap N(y)|$ denote the number of common
neighbours of $x$ and $y$, and let $d_x = \deg(x)$, $d_y = \deg(y)$.
Then:
\begin{equation}
\kappa(x,y) \leq \frac{t_{xy}}{\min(d_x, d_y)}.
\label{eq:orc_upper}
\end{equation}
\end{theorem}

\begin{proof}
The optimal coupling between $\mu_x$ and $\mu_y$ matches each
common neighbour $z \in N(x) \cap N(y)$ at zero transport cost.
The minimum coupling assigns at most $t_{xy}/\min(d_x, d_y)$ of the
total mass at zero cost, so
$W_1(\mu_x, \mu_y) \geq 1 - t_{xy}/\min(d_x, d_y)$,
giving $\kappa(x,y) \leq t_{xy}/\min(d_x, d_y)$.
\end{proof}

\begin{remark}[On the Lower Bound]
\label{rem:lower_bound}
It is natural to conjecture a matching lower bound
$\kappa(x,y) \geq t_{xy}/\max(d_x,d_y) - |d_x-d_y|/\max(d_x,d_y)$.
We report that it does not hold in general. Numerical experiments with
exact Wasserstein-1 computation on Erdős-Rényi and Barabási-Albert
networks show it fails in networks with sparse triangle structure: when
$t_{xy} = 0$ on most edges, many non-common neighbours lie
at geodesic distance $> 1$, making the actual transport cost
higher than the bound assumes. The correct lower bound in full
generality requires bounding the geodesic distances between
non-common neighbours, a condition that holds for
triangle-rich networks (clustering coefficient $C > 0.3$)
but fails for sparse random graphs. Establishing tight lower
bounds for arbitrary graphs is left as an open problem.
\end{remark}

\begin{proposition}[Non-negativity Condition for Lin-Lu-Yau Curvature]
\label{prop:orc_nonneg}
Let $\{x,y\} \in E(G)$ with $d_x = \deg(x) < d_y = \deg(y)$, let
$\Delta = N(x) \cap N(y)$ and $\delta = |\Delta|$, and write
$A = N(x)\setminus N[y]$ and $B = N(y)\setminus N[x]$.
By~\cite{lly2026nonreg} (Theorem~1.2), if
\begin{equation}
\delta \geq \frac{2d_x d_y - 2d_x - 2d_y}{d_y + 2d_x},
\label{eq:lly_nonneg}
\end{equation}
then the Lin-Lu-Yau Ricci curvature~\cite{lin2011}
$\kappa^{\mathrm{LLY}}(x,y) \geq 0$.
Moreover, $\kappa^{\mathrm{LLY}}(x,y) > 0$ unless simultaneously:
(i) equality holds in~\eqref{eq:lly_nonneg},
(ii) there are no edges between $\Delta$ and $B$, and
(iii) $\mathrm{dist}(u,w) = 3$ for all $u \in A$, $w \in B$.

\textbf{Remark.} Condition~\eqref{eq:lly_nonneg} makes non-negative curvature
sufficient once the common neighbourhood is \emph{large} relative to the
degrees: the right-hand side grows like $2d_x/3$ when $d_y$ is much the larger
of the two, so the condition asks the edge to close a number of triangles
comparable to the degree of its lower-degree endpoint. The hypothesis
$d_x < d_y$ is strict in~\cite{lly2026nonreg}, so the statement does not
specialise to regular graphs, for which the exact formulas of that work
(Theorem~1.1) apply instead.

The general bounds on $\kappa^\alpha(x,y)$ for all $\alpha \in [0,1]$ and
arbitrary degree sequences are given by Theorem~5.3 of~\cite{lly2026nonreg},
which extends earlier results of Jost-Liu~\cite{jost2014} to non-regular graphs.

\textbf{Connection to Tr$(A^3)$.} The condition~\eqref{eq:lly_nonneg}
depends on $\delta = t_{xy}$, the same common-neighbour count that
determines the local contribution to $\operatorname{Tr}(A^3) =
\sum_{(x,y)\in E} 2 t_{xy}$, and it points the same way as
Theorem~\ref{thm:orc_upper}: triangle-rich edges are the ones whose curvature
is non-negative. Edges satisfying~\eqref{eq:lly_nonneg}
contribute non-negatively to the total Ollivier-Ricci curvature,
while edges with $\delta < \frac{2d_x d_y - 2d_x - 2d_y}{d_y + 2d_x}$
are the ones that may have $\kappa < 0$ and act as topological bridges. Those
are the edges that close few triangles and therefore contribute little to
$\operatorname{Tr}(A^3)$, so the curvature reading and the triadic one agree on
which edges are which. What they do \emph{not} agree on is the direction of the
composite: a triangle-poor bridge is a negative-curvature edge, and it is the
distribution of the triangles across nodes, read by $\mathrm{Coex}$ rather than
by $\operatorname{Tr}(A^3)$ alone, that carries the monopoly signature
(Remark~\ref{rem:direction}).
\end{proposition}

\begin{corollary}[Threshold for Positive Total Curvature]
\label{cor:positive_total}
The upper bound $\Sigma\kappa^+ \leq \operatorname{Tr}(A^3)/(2d_{\min})$
holds on every network in Table~\ref{tab:orc_bounds}. It is an
inequality and not an estimate: edges satisfying~\eqref{eq:lly_nonneg}
contribute non-negatively to $\Sigma\kappa^+$, whereas edges with
$\kappa < 0$ contribute nothing to it while their triangles still
enter $\operatorname{Tr}(A^3)$. The bound is therefore expected to be
loose in graphs where a large fraction of edges act as bridges, and
looser still when $d_{\min}$ is small; see
Remark~\ref{rem:loose_bound}.
\end{corollary}

\textbf{Consequence for Tr$(A^3)$.}
Since $\text{Tr}(A^3) = 6N_\triangle = \sum_{(x,y)\in E} 2\,t_{xy}$,
the upper bound gives:
\begin{equation}
\sum_{(x,y) \in E} \kappa^+(x,y) \leq
\frac{\text{Tr}(A^3)}{2\,d_{\min}}.
\label{eq:tra3_upper}
\end{equation}

\textbf{Empirical validation.}
Using exact $W_1$ computation via optimal transport (Wasserstein-1
distance between neighbour distributions), we validate on six
networks. Results in Table~\ref{tab:orc_bounds} show that all six
satisfy the upper bound~\eqref{eq:tra3_upper}. The bound is, however,
far from tight, and its slack does not order by clustering: the ratio
$\Sigma\kappa^+/\text{UB}$ is $0.065$ for Karate Club, essentially the
same as $0.047$ for the Erd\H{o}s-R\'{e}nyi graph and $0.069$ for the
Barab\'{a}si-Albert graph, even though the first has a dense triangle
structure and the other two do not. Remark~\ref{rem:loose_bound}
identifies what does control the slack. The companion lower bound
$\operatorname{Tr}(A^3)/(2d_{\max})$, which is natural to conjecture
alongside~\eqref{eq:tra3_upper}, is not satisfied on the two sparse
random graphs: it evaluates to $9.00$ against $\Sigma\kappa^+ = 1.90$
for Erd\H{o}s-R\'{e}nyi and to $7.83$ against $3.25$ for
Barab\'{a}si-Albert. This is the aggregate counterpart of the
edge-wise failure reported in Remark~\ref{rem:lower_bound}, and it is
why only the upper bound is claimed here.

\begin{table}[ht]
\centering
\caption{ORC upper bound validation. $\Sigma\kappa^+$ is computed with exact
Wasserstein-1 optimal transport; LB $= \text{Tr}(A^3)/(2d_{\max})$ is the
lower bound one would naturally conjecture alongside it, and
UB $= \text{Tr}(A^3)/(2d_{\min})$ is the bound of~\eqref{eq:tra3_upper}.
All six networks satisfy the upper bound; the conjectured lower bound fails
on the two sparse random graphs. Slack is governed by $d_{\min}$, not by
clustering; see Remark~\ref{rem:loose_bound}. The values were computed with
\texttt{colabs/fsri\_colab4b\_orc\_exact.py}, which is included in the
accompanying repository.}
\label{tab:orc_bounds}
\small
\begin{tabular}{lrrrrrrr}
\toprule
Network & $N$ & $M$ & $d_{\min}$ & $d_{\max}$ & LB & $\Sigma\kappa^+$ & UB \\
\midrule
Karate Club~\cite{zachary1977}         & 34 & 78  & 1  & 17 & 7.94   & 8.74   & 135.00  \\
Florentine Families~\cite{padgett1993} & 15 & 20  & 1  & 6  & 1.50   & 1.50   & 9.00    \\
Les Mis\'{e}rables                     & 77 & 254 & 1  & 36 & 38.92  & 82.40  & 1401.00 \\
Davis Southern Women                   & 18 & 139 & 11 & 17 & 111.35 & 111.78 & 172.09  \\
\midrule
Erd\H{o}s-R\'{e}nyi ($n{=}40$)         & 40 & 110 & 2  & 9  & 9.00   & 1.90   & 40.50   \\
Barab\'{a}si-Albert ($n{=}40$)         & 40 & 111 & 3  & 18 & 7.83   & 3.25   & 47.00   \\
\bottomrule
\end{tabular}
\end{table}

\begin{remark}[Why the Bound Is Loose]
\label{rem:loose_bound}
The slack in~\eqref{eq:tra3_upper} is governed by $d_{\min}$ rather
than by clustering. A single pendant vertex sets $d_{\min} = 1$ and
doubles the bound relative to an otherwise identical graph with
$d_{\min} = 2$, however many triangles the graph contains. Three of
the six networks in Table~\ref{tab:orc_bounds} have $d_{\min} = 1$
(Karate Club, Florentine Families, Les Mis\'{e}rables), and for these
the reported UB is exactly $\text{Tr}(A^3)/2$: $135.00 = 270/2$,
$9.00 = 18/2$ and $1401.00 = 2802/2$, for $45$, $3$ and $467$
triangles respectively. This is why Karate Club and Les
Mis\'{e}rables, both triangle-rich, show slack indistinguishable from
that of the sparse random graphs, while Davis Southern Women, the only
network here with $d_{\min} = 11$, is the one case where the bound is
nearly attained ($111.78$ against $172.09$). A second and independent
effect acts in the sparse random graphs, where most edges have
$\kappa < 0$ and act as bridges rather than closing triangles, so that
$\Sigma\kappa^+$ is small in absolute terms as well. Neither effect is
evidence against the bound, and neither is captured by the clustering
coefficient.
\end{remark}

This gives $\mathcal{T}(W)$ a precise geometric meaning: it measures
the \emph{total positive curvature} of the weight graph. Minimising
$\mathcal{T}(W)$ flattens the curvature distribution, moving the
network toward a configuration where curvature is uniformly
distributed, the network analogue of a maximally symmetric
Riemannian manifold.

\begin{remark}[Continuum limits and other curvature observables: what is
not claimed here]
\label{rem:continuum}
Two results in the literature extend the discrete construction toward
continuous geometry, and we record them as context rather than as
contributions of our own. The Ollivier-Ricci curvature of random geometric
graphs embedded in a Riemannian manifold converges in the continuum limit to
the Ricci curvature of that manifold, and the result extends to Lorentzian
signature~\cite{vanderhoorn2021}; building on it, the construction has been
defined for causal sets using Lorentzian optimal
transport~\cite{orc_causalsets2026}. Where a network encodes causal relations
in that sense~\cite{bombelli1987}, its curvature approximates a Ricci scalar,
and one may read the two ends of $\Omega$ as regions where geodesics focus
and defocus. We derive nothing in that direction, and a reader who discounts
this remark entirely loses nothing that the rest of the paper relies on.
Separately, nonperturbative quantum gravity has its own curvature observable
for discrete geometries~\cite{klitgaard2018a,klitgaard2018b,klitgaard2020},
which is a different construction: it compares distances between spheres
rather than transporting neighbour measures, and admits no triangle-counting
identity, so the bound above does not transfer to it. Whether $\mathcal{T}$
has an analogue for quantum channels under non-commutative transport
metrics~\cite{gao2021} is open.
\end{remark}
\section{Application: Omega-S Neural Network Regulariser}
\label{sec:omega_s}
% ============================================================

\subsection{From Network Theory to Neural Networks}

The FSRI framework has an immediate application to neural network
regularisation. For a weight matrix $W^{(\ell)} \in \mathbb{R}^{m \times n}$
in layer $\ell$ of a neural network, we define the layer-wise
topological pressure:
\begin{equation}
\mathcal{T}_\ell = \frac{\lambda_\Omega}{|\operatorname{numel}(W^{(\ell)})|}
\cdot \widehat{\operatorname{Tr}}(W^{(\ell)}W^{(\ell)\top})^3,
\label{eq:omega_layer}
\end{equation}
and the total regularisation penalty:
\begin{equation}
\mathcal{L}_\mathcal{T} = \sum_\ell \mathcal{T}_\ell.
\label{eq:total_penalty}
\end{equation}

The Omega-S regulariser adds $\mathcal{L}_\mathcal{T}$ to the task loss every
$K$ training steps. By Proposition~\ref{prop:gradient_flow}, this drives the network's weight
matrices toward minimum-entropy-production configurations, that is, states
where structural monopolies are suppressed.

One point should be stated plainly, because the shared name invites the
opposite reading. \emph{The regulariser minimises $\mathcal{T}(W)$ and never
computes $\Omega$.} No clustering coefficient, connection density,
modularity or degree variance enters the training loop: the penalty is the
Hutchinson estimate of $\operatorname{Tr}(WW^\top)^3$ and nothing else. The
index and the penalty are related by interpretation, through
Definition~\ref{def:tep} and Proposition~\ref{prop:gradient_flow}, and not
by computation. This is why the tension identified in Remark~\ref{rem:direction}
does not arise in the applied setting: what is minimised is the topological
pressure, whose thermodynamic reading is unambiguous, rather than the
composite index, whose response to that minimisation depends on how
$\mathrm{Coex}$ moves alongside it and is therefore a measured outcome.

\subsection{Empirical Validation}

The Omega-S regulariser~\cite{acedo2026omegas} validates the theoretical
predictions:

\begin{table}[ht]
\centering
\caption{Omega-S empirical results, with the setting for each row. All
values are as reported in the applied paper~\cite{acedo2026omegas}.
Retention is measured on Llama-3-8B (LoRA, code $\to$ prose) as
HumanEval pass@1 post-prose / post-code; the degree-variance row is measured
on the three-layer MLP with Split-MNIST and the structured-sparsity row on
GPT-2, not on Llama-3-8B.}
\label{tab:omega_results}
\small\setlength{\tabcolsep}{4pt}
\begin{tabular}{lcc}
\toprule
Result & Value & Interpretation \\
\midrule
Code capability kept & $0.173 \to 0.238$, 9/10 &
  Primary metric, vs no regulariser; \\
 & & sign test $p=0.011$ \\
Same, as retention ratio & $62.9\% \to 84.1\%$ &
  The same comparison, in ratio form \\
Retention vs weight decay & $61.9\% \to 84.1\%$, 10/10 &
  Secondary, $p=0.002$; thin selection \\
Retention vs EWC & $69.1\% \to 84.1\%$, 8/10 &
  Secondary, $p=0.014$; thin selection \\
vs row-norm control (own sweep) & 8/10, $p=0.055$ &
  Not reducible to norm balancing \\
Degree variance reduction (MLP) & $41.94 \to 0.136$ &
  $8\times$ the reduction obtained by \\
 & & weight decay ($1.12$) \\
Structured sparsity (GPT-2) & $-0.72\%$ vs $-0.77\%$ &
  No gain over group lasso alone; the \\
 & & compression is group lasso's \\
Training overhead & $<4\%$ per step &
  Analytic, penalty every 10 steps \\
Latency overhead (2-GPU FSDP) & $+1.5\%$ &
  Measured; zero inter-GPU cost \\
\bottomrule
\end{tabular}
\end{table}

The retention result is particularly significant from the thermodynamic
perspective: minimum-entropy-production states (low $\mathcal{T}$) are
more stable under perturbation (new training data), consistent with
Prigogine's characterisation of stationary non-equilibrium states as
maximally robust against small deviations.

% ============================================================
\section{Applications Beyond Neural Networks}
\label{sec:applications}
% ============================================================

The FSRI framework is domain-agnostic, and the readings below are offered
as interpretations of the same functional in other settings rather than as
results. None of them is measured here.

\textbf{Ecological networks.} A network dominated by a single node
($\mathrm{Coex} \to \infty$, $C \to 0$) collapses $\Omega \to 0$, which the
framework reads as systemic fragility; the Florentine Medici marriage and
credit network~\cite{padgett1993} is the standard illustration of that
topology. Soil microbial networks under resource abundance exhibit the same
monopoly structure~\cite{ortiz2021}, and land-use intensification is reported
to reduce network complexity and the abundance of keystone
taxa~\cite{banerjee2019}. The one experimental test of the direction this
reading assumes manipulates soil biodiversity in grassland microcosms and
finds that network complexity predicts ecosystem multifunctionality, with the
relation surviving a correction for richness~\cite{wagg2019}; that result is
about complexity and function rather than about $\Omega$, and we cite it as
the motivation for the numerator of~\eqref{eq:fsri} and not as a validation of
the composite.

\textbf{Financial systems.} The Minsky Moment~\cite{minsky1986}, the
sudden collapse of a leveraged financial system, is a structural
phase transition from high-$\Omega$ (distributed, mutualistic) to
low-$\Omega$ (monopolistic, fragile) topology under energy abundance
($E_\text{ext}$ = credit availability).

\textbf{Power grids.} Cascade failure in a linear (high-$M$, low-$C$) grid
reads in the same terms, as a topological phase transition triggered by a
central shock.

\textbf{Medical systems.} Microbiome dysbiosis is sometimes described in the
same terms, with energy abundance driving specialisation ($M$ increases, $C$
decreases, $\mathrm{Coex}$ increases, $\Omega$ collapses). We flag this
reading as the least supported of those listed here, and as an illustration
rather than an expectation. Evidence from the human gut does not line up with
it uniformly: co-occurrence networks in inflammatory bowel disease are
reported to be reorganised rather than simply degraded, and the healthy
network is the one more easily disrupted by removal of a few key
taxa~\cite{baldassano2016}. The orientation of $M$ in this setting is subject
to Remark~\ref{rem:modularity_scope}. The reading is also subject
to the sampling condition of Section~\ref{sec:limitations}: the node set of
a co-occurrence network of taxa grows with sequencing depth, so readings
taken at unequal depth are not comparable without an explicit
standardisation, and the factor most affected, $\mathrm{Coex}$, is the one
this reading relies on.

\textbf{Multi-agent language-model systems.} A recent analysis of communication
graphs among language-model agents derives a two-sided constraint on
connectivity of exactly the shape this framework describes, and reaches it from
epidemiology rather than from triadic structure~\cite{niu2026}. A clean-task
reliability target imposes a minimum degree, because a node needs enough
independent opinions for a majority vote to reach the required accuracy; an
error-propagation target imposes a maximum, because at fixed exposure per edge
the invasion factor grows with the spectral radius of the adjacency matrix.
The feasible regular degrees are then confined to an interval, and the authors
characterise when that interval is empty. The construction differs from ours in
every detail, and we record the structural agreement rather than any numerical
one; its relation to Definition~\ref{def:fsri} is the subject of
Remark~\ref{rem:window}.

% ============================================================
\section{Limitations and Future Work}
\label{sec:limitations}
% ============================================================

We set out the scope of each result, distinguishing what is established
from what remains conjectural.

\textbf{Scope of the ORC-triangle connection.} Proposition~\ref{prop:orc_triangles} is proved
using Jost-Liu~\cite{jost2014} for regular and semi-regular graphs. The
approximation $\kappa^+(x,y) \propto |N(x) \cap N(y)|/\bar{d}$ fails for
highly heterogeneous degree distributions; a fully general proof requires
the exact Wasserstein-1 computation and is left for future work.

\textbf{Scope of the FSRI proportionality.} Proposition~\ref{prop:ws_exact} derives the
exact FSRI formula for Watts-Strogatz networks,
$\Omega = [(3K-6)/(4K-2)] \cdot (1-p)^2/(np)$, which shows that the
proportionality $\Omega \propto (S_\text{env}/E_\text{ext})^2$ does not hold
in general. Substituting $\mathrm{Var}(k) = Kp(1-p)$ and $E_\text{ext} = K$
into~\eqref{eq:omega_ws} gives $(S_\text{env}/E_\text{ext})^2 = p(1-p)/K$ and
hence the exact identity
\[
\Omega \;=\; \frac{1}{nK}\cdot\frac{3K-6}{4K-2}\cdot
\frac{(1-p)^3}{(S_\text{env}/E_\text{ext})^{2}},
\]
in which the ratio appears in the \emph{denominator}. The relation is inverse
rather than proportional, which is what produces the negative correlations
$r = -0.407$ to $-0.573$ reported above rather than merely weak positive ones.
Deriving the corresponding expression for other network models
(Barabási-Albert, Erdős-Rényi) remains open.

\textbf{Sampled graphs and the growing node set.}
Remark~\ref{rem:comparability} states the condition under which two values
of $\Omega$ may be compared: matched order and matched density. A stronger
condition applies when $A$ is not given but estimated from a finite sample
whose node set grows with the sampling effort. Co-occurrence networks of
microbial taxa at different sequencing depths, and lexical or syntactic
networks at different corpus sizes, are of this kind: additional sampling
adds vertices rather than only refining the weights among those already
present. In ecological terms this is the sampling regime of a metacommunity
with a non-zero speciation rate~\cite{hubbell2001}, in which the local
sample is drawn from a regional pool that is itself still being added to,
so that richness does not saturate and neither does any factor computed
from it. Matching the sample size is then not sufficient, because two units
of equal sample size sit at different points of their respective
accumulation curves.

We measured the effect on syntactic dependency networks built from
Universal Dependencies treebanks for 28 languages, at token budgets from
$2\,000$ to $32\,000$, fitting $\log_{10}(\cdot) = a + b \log_{10} N$ per
language. The accumulation exponents are $\beta_V = 0.77 \pm 0.06$ for the
vertex set and $\beta_E = 0.955 \pm 0.017$ for the edge set. The former is
the sublinear vocabulary growth of Herdan and Heaps~\cite{heaps1978},
measured here at the small end of the corpus range where that exponent is
known to exceed the $0.4$ to $0.6$ typical of large corpora. The latter
says that distinct dependency pairs accumulate almost linearly with the
sample: the edge set is nowhere near saturation in this range, and the
tightness of $\beta_E$ across typologically distant families is worth noting
in itself.

All four factors drift, and the difference between them is one of degree
rather than of kind: $b = +1.03$ for $\mathrm{Coex}$ and $-0.59$ for $D$,
against $-0.14$ for $C$ and $-0.09$ for $M$. For the index itself
$b = -1.68$, so a tenfold increase in corpus size divides $\Omega$ by
roughly fifty. Two of these slopes carry different weight, and we separate
them rather than report them alike. That of $D$ follows algebraically from
$\beta_V$ and $\beta_E$, since $D$ is a fixed function of $n$ and $m$, and is
therefore a consistency check on the fits and not independent evidence.
That of $\mathrm{Coex}$ does not, since it depends on the shape of the
degree distribution and not only on the two accumulation rates, and it is
the slope that carries information. The drift is not an artefact of the
small end of the range: the slope of $\mathrm{Coex}$ changes by $-0.02$
between the lower and upper halves of the budgets tested. Neither does it
extrapolate, since $\beta_V$ is itself known to fall as corpora
grow~\cite{gerlach2013}.

Two routes are available in the literature. The first is to standardise by
sample \emph{coverage} rather than by sample size, following Chao and
Jost~\cite{chao2012coverage}, who show that equal-sized samples
systematically bias comparisons because a size sufficient to characterise a
poorer community is insufficient for a richer one, and that coverage is
estimable from the singleton count. The second is to replace the
concentration factor by a parameter of the generating process rather than a
statistic of the realised sample, for instance the fitted tail exponent
$\gamma$ of the degree distribution~\cite{clauset2009}. On these corpora
that exponent is estimable and lies below three throughout
($\gamma = 2.67 \pm 0.22$ over 582 samples, range $1.92$ to $2.99$), which
makes the divergence condition of Remark~\ref{rem:comparability} concrete
rather than hypothetical. It does not account for the observed rate: the
standard scaling for a tail exponent of that size predicts a drift about an
order of magnitude smaller than the one measured. It also changes what the
factor means, breaking the reading of $\mathrm{Coex}$ as a degree variance
on which the competitive-exclusion analogy of Remark~1 rests.

The same two factors identified analytically in
Proposition~\ref{prop:shape_invariant} are the ones that carry the drift
here, and the shape-normalised index of Definition~\ref{def:shape} largely
removes it. The Gini coefficient of the degree sequence drifts at
$b = +0.09$ and the normalised entropy of that sequence at $b = -0.02$,
against $+1.03$ for $\mathrm{Var}(\mathbf{k})$. Substituting the
concentration factor alone leaves the index at $b = -0.73$, because $D$
still drifts at $-0.59$ and, as noted above, the slope of $D$ is
$\beta_E - 2\beta_V$, a property of the sampling regime rather than of the
statistic chosen. Replacing $D$ by the mean degree as well, whose slope is
$\beta_E - \beta_V = 0.19$, gives $b = +0.04$ for $\widetilde{\Omega}$, a
factor of roughly $1.1$ per decade of corpus size against the factor of
$47$ carried by $\Omega$ as defined. Restoring comparability therefore
requires changing two factors rather than one, and the pair that works
empirically is the pair the proposition identifies.

We none the less keep $\Omega$ as the index of Definition~\ref{def:fsri}
and $\widetilde{\Omega}$ as a variant for sampled graphs, for the reason
given in Remark~\ref{rem:shape_scope}: the proposition covers graphs with
homogeneous degrees, the measurement covers estimated graphs with
heavy-tailed ones, and the second does not follow from the first. What can
be said is that the variant is not an ad hoc repair. It has the same
analytical status on the Watts-Strogatz family as the index it modifies,
and it is the substitution that the scale argument predicts in advance.

Two elements of the treatment above have published precedents that are worth
naming, because both are sometimes taken to be idiosyncratic. The first is the
construction: inferring a single association network over all units and then
reading, for each unit, the subnetwork induced by the taxa it contains is the
procedure used to relate microbiome complexity to ecosystem multifunctionality
in grassland microcosms~\cite{wagg2019}, and a per-sample connectivity
statistic computed from a global matrix of pairwise associations is the basis
of the cohesion metric of~\cite{herren2017}. The second is the correction:
the same work observes that network complexity depends on the number of nodes,
detrends for richness before interpreting, and reports that the residual
relation to function survives that detrending~\cite{wagg2019}. The condition
stated here is therefore not a new objection to an established practice but
the quantitative form of a correction that practice already applies.

A limitation applies to the measurement itself. Universal Dependencies
provides several treebanks of differing genre for some languages, and the
dispersion of these quantities between treebanks of one language is not
small relative to the dispersion between languages. Measured on sixteen
treebanks across five languages at a matched budget, the ratio of the
second to the first is $3.6$ for the Gini coefficient, $3.2$ for the
normalised entropy and $2.15$ for $\mathrm{Var}(\mathbf{k})$. Comparison of
network quantities estimated from treebanks is therefore genre-contaminated
across languages as well as sampling-dependent within them, and the margins
above should be read accordingly.

\textbf{Status of the ORC numerical validation.} The Colab~3 experiment
confirms numerically that ORC $\kappa(x,y)$ correlates strongly with edge
contribution to $\operatorname{Tr}(A^3)$ in Watts-Strogatz networks,
supporting Proposition~\ref{prop:orc_triangles} empirically. The evidence
is numerical rather than analytical, and restricted to that family.

\textbf{Scope of the spectral identity.}
Proposition~\ref{prop:spectral_identity} constrains what can be claimed for
the global index: as a scalar detector, $\mathcal{T}$ should be benchmarked
against, and is expected to tie, other spectral functionals. Claims of
superiority over spectral methods are defensible only for the node-wise
field (Corollary~\ref{cor:locus}) or for the training-time gradient flow,
not for the global statistic. The signed extension is stated but not yet
validated empirically; whether signed $\mathcal{T}$ improves detection in
domains where sign is meaningful is an open question.

The geometric reading is limited to what Section~\ref{sec:orc_graphs}
establishes, namely the inequality relating common-neighbour counts to
$\operatorname{Tr}(A^3)$ and its measured slack. Continuum limits, causal
sets and the curvature observables of nonperturbative quantum gravity are
noted in Remark~\ref{rem:continuum} as context; we derive no result in that
direction and none of our claims depends on one.

% ============================================================
\section{Conclusion}
% ============================================================

We have presented a mathematical framework connecting the Functional
Symbiotic Resilience Index $\Omega$ to non-equilibrium thermodynamics
via the topological pressure functional $\mathcal{T}(W) =
\operatorname{Tr}(WW^\top)^3$.

The key results are:
\begin{enumerate}
\item $\mathcal{T}(W)$ is proportional to the global clustering
coefficient of the induced weight graph, providing a precise
topological interpretation.
\item The gradient $3(WW^\top)^2$ is the \emph{topological pressure}
on each weight, whose Gini coefficient discriminates monopoly networks
($G = 0.986$) from uniform networks ($G = 0.272$) by 3.6×.
\item The gradient flow $\dot{W} = -\lambda \nabla_W \mathcal{T}(W)$
implements Prigogine's minimum-entropy-production principle in the
space of weight configurations.
\item Phase transitions in triangle density (fugacity mechanism)
confirm the thermodynamic analogy: Tr$(A^3)$ increases 76× over the
transition, while the Ollivier-Ricci curvature interpretation provides
a geometric framework consistent with the relational motivation of the
index.
\item The index is not invariant to the order of the graph, and the
dependence is carried by exactly two of its four factors. Replacing them
by their shape-normalised counterparts, the mean degree and the Gini
coefficient of the degree sequence, removes the dependence exactly on
Watts-Strogatz graphs (Proposition~\ref{prop:shape_invariant}) and reduces
it from a factor of $47$ to a factor of $1.1$ per decade of sampling effort
on estimated graphs (Section~\ref{sec:limitations}).
\end{enumerate}

Since several of the results above are restrictive, we state in one place
what the framework claims after them. Three things survive. The node-wise
field $\operatorname{diag}(A^3)$ is not a function of the spectrum
(Proposition~\ref{prop:node_nonspectral}) and is therefore the component
that no eigenvalue-based scalar can reproduce. The training-time gradient
flow of Proposition~\ref{prop:gradient_flow} is a construction rather than a
statistic, and its empirical outcome is reported in
Section~\ref{sec:omega_s}. And the global scalar $\Omega$ remains a usable
descriptive index within the conditions of use set out in
Remark~\ref{rem:comparability} and Section~\ref{sec:limitations}, namely at
matched graph order and density and, for graphs estimated from samples,
under an explicit standardisation, or in the shape-normalised form of
Definition~\ref{def:shape} where those conditions cannot be met. What we do
not claim is that the global
scalar dominates other spectral functionals as a detector; Proposition~\ref{prop:spectral_identity}
forbids it, and we say so rather than leave it to be found.

The Omega-S neural network regulariser is a concrete instantiation of the
framework rather than a proof of it. Minimising $\mathcal{T}(W)$ during
training yields networks that are more stable under sequential learning
(mean retention $61.9\%\to84.1\%$ over tuned weight decay across ten seeds, ten
of ten positive) and, in combination with group lasso, structurally prunable,
although the compression obtained in that setting is attributable to the
group-lasso term and the marginal contribution of the penalty remains to be
isolated. These are empirical outcomes consistent with the
minimum-entropy-production reading developed above; they do not establish
that reading, which would require measuring entropy production directly
rather than inferring it from the structural proxy. We record the
distinction because the thermodynamic interpretation is the more
speculative part of this work, and the empirical results stand
independently of whether it is ultimately upheld.

% ============================================================
\bibliographystyle{unsrt}
% ============================================================
\section*{Declaration of LLM Usage}
% ============================================================

Large language models were used as writing and coding assistants during the
preparation of this manuscript: to improve language clarity and
presentation, to assist with coding and debugging for the numerical
verifications, and to help identify related work. All technical ideas,
theoretical results, proofs, experimental design and final interpretations
were developed, verified and approved by the author.

\begin{thebibliography}{99}

\bibitem{acedo2026omegas}
Acedo A. (2026).
Omega-S: A Functional Resilience Index for LLM Fine-Tuning. Less Forgetting
Without Past Data, and More Reliable Structured Pruning.
\textit{arXiv}:2608.03887. \url{https://arxiv.org/abs/2608.03887}

\bibitem{acedo2022indices}
Acedo A, Ortega-Arranz H, Almonacid D, Ferrero A (2022).
Methods and systems for generating and applying agronomic indices from
microbiome-derived parameters.
US Patent Application Publication No.\ US~2022/0268756~A1
(Appl.\ No.\ 17/665,332, filed 4 February 2022), Biome Makers Inc.

\bibitem{biome_patent}
Ortiz-Alvarez R, Ferrero A, Imam N, Acedo A, Belda-Aguilar I, Almonacid D
(2021).
Methods and systems for evaluating ecological disturbance of an agricultural
microbiome based upon network properties of organism communities.
WIPO Publication No.\ WO~2021/119528~A1 (PCT/US2020/064668), Biome Makers Inc.

\bibitem{prigogine1967}
Prigogine I. (1967).
\textit{Introduction to Thermodynamics of Irreversible Processes}.
Wiley-Interscience.

\bibitem{watts1998}
Watts D J, Strogatz S H. (1998).
Collective dynamics of `small-world' networks.
\textit{Nature}, 393, 440--442.

\bibitem{barabasi1999}
Barabási A-L, Albert R. (1999).
Emergence of scaling in random networks.
\textit{Science}, 286(5439), 509--512.

\bibitem{hutchinson1990}
Hutchinson M F. (1990).
A stochastic estimator of the trace of the influence matrix for
Laplacian smoothing splines.
\textit{Communications in Statistics --- Simulation and Computation},
19(2), 433--450.

\bibitem{ollivier2009}
Ollivier Y. (2009).
Ricci curvature of Markov chains on metric spaces.
\textit{Journal of Functional Analysis}, 256(3), 810--864.

\bibitem{motif2026}
Anonymous. (2026).
Motif enrichment as a driver of scale-free behavior in rewired random
regular graphs.
\textit{arXiv:2604.23152}.

\bibitem{martins2023}
Martins G F, et al. (2023).
Topologically-constrained fluctuations and thermodynamics regulate
nonequilibrium response.
\textit{arXiv:2305.19348}.

\bibitem{toporesilience2017}
Asllani M, et al. (2017).
Topological resilience in non-normal networked systems.
\textit{arXiv:1706.02703}.

\bibitem{strauss1986}
Strauss D. (1986).
On a general class of models for interaction.
\textit{SIAM Review}, 28(4), 513--527.

\bibitem{lu2024alpha}
Lu Y, et al. (2024).
AlphaPruning: Using Heavy-Tailed Self Regularization Theory for
Improved Layer-wise Pruning of Large Language Models.
\textit{arXiv:2410.10912}.

\bibitem{ortiz2021}
Ortiz-Álvarez R, et al. (2021).
Consistent Changes in the Soil Bacterial Communities of Agricultural
Ecosystems.
\textit{mSystems}, 6(1).

\bibitem{minsky1986}
Minsky H. (1986).
\textit{Stabilizing an Unstable Economy}.
Yale University Press.

\bibitem{zachary1977}
Zachary W W. (1977).
An information flow model for conflict and fission in small groups.
\textit{Journal of Anthropological Research}, 33(4), 452--473.

\bibitem{padgett1993}
Padgett J F, Ansell C K. (1993).
Robust Action and the Rise of the Medici, 1400--1434.
\textit{American Journal of Sociology}, 98(6), 1259--1319.

\bibitem{lly2026nonreg}
Jung J, Kwon Y S, Lee S, Park J. (2026).
A combinatorial approach to $\alpha$-Ricci and Lin-Lu-Yau Ricci
curvatures on graphs.
\textit{arXiv:2606.24261v1} [math.DG], 23 June 2026.
Global Basic Research Laboratory, DGIST, Republic of Korea.

\bibitem{lin2011}
Lin Y, Lu L, Yau S-T. (2011).
Ricci curvature of graphs.
\textit{Tohoku Mathematical Journal}, 63(4), 605--627.

\bibitem{barrat2000}
Barrat A, Weigt M. (2000).
On the properties of small-world network models.
\textit{European Physical Journal B}, 13(3), 547--560.

\bibitem{sandhu2016}
Sandhu R S, Georgiou T T, Tannenbaum A R. (2016).
Ricci curvature: an economic indicator for market fragility and systemic risk.
\textit{Science Advances}, 2(5), e1501495.

\bibitem{expost2024}
Anonymous. (2024).
On the Ollivier-Ricci curvature as fragility indicator of the stock markets.
arXiv:2405.07134.

\bibitem{bauer2012}
Bauer F, Jost J, Liu S. (2012).
Ollivier-Ricci curvature and the spectrum of the normalized graph Laplace
operator.
\textit{Mathematical Research Letters}, 19(6), 1185--1205.

\bibitem{jost2014}
Jost J, Liu S. (2014).
Ollivier's Ricci curvature, local clustering and curvature-dimension
inequalities on graphs.
\textit{Discrete \& Computational Geometry}, 51(2), 300--322.

\bibitem{schnakenberg_nets}
Maes C, Netočný K. (2007).
Minimum entropy production principle from a dynamical fluctuation law.
\textit{Journal of Mathematical Physics}, 48, 053306.

\bibitem{prigogine_critical}
Bertola M, Grunbaum F A. (2007).
A critical analysis of the minimum entropy production theorem and its
application to heat and fluid flow.
\textit{International Journal of Heat and Mass Transfer}, 51(5--6), 1907--1912.

\bibitem{maxcaliber2018}
Ghosh K, Dixit P D, Agozzino L, Dill K A. (2018).
Maximum caliber: a general variational principle for dynamical systems.
\textit{arXiv:1505.05479}.

\bibitem{samal2021}
Samal A, Pharasi H K, Ramaia S J, Kannan H, Saucan E, Jost J, Chakraborti A.
(2021).
Network geometry and market instability.
\textit{Royal Society Open Science}, 8, 201734.

\bibitem{sia2019}
Sia J, Jonckheere E, Bogdan P. (2019).
Ollivier-Ricci curvature-based method to community detection in complex networks.
\textit{Scientific Reports}, 9, 9800.

\bibitem{vanderhoorn2021}
van der Hoorn P, Cunningham W J, Meskhidze H, Krioukov D. (2021).
Ollivier-Ricci curvature convergence in random geometric graphs.
\textit{Physical Review Research}, 3, 013211.

\bibitem{orc_causalsets2026}
Anonymous. (2026).
Ollivier-Ricci Curvature for Causal Sets.
\textit{arXiv:2606.04910}.

\bibitem{bombelli1987}
Bombelli L, Lee J, Meyer D, Sorkin R D. (1987).
Space-time as a causal set.
\textit{Physical Review Letters}, 59(5), 521--524.

\bibitem{cartwright1956}
Cartwright D, Harary F (1956).
Structural balance: a generalization of Heider's theory.
\emph{Psychological Review} 63(5):277--293.

\bibitem{klein1954}
Klein M J. (1954).
Principle of minimum entropy production.
\textit{Physical Review}, 96, 250--255.

\bibitem{klitgaard2018a}
Klitgaard N, Loll R. (2018).
Introducing quantum Ricci curvature.
\textit{Physical Review D}, 97(4), 046008. \textit{arXiv:1712.08847}.

\bibitem{klitgaard2018b}
Klitgaard N, Loll R. (2018).
Implementing quantum Ricci curvature.
\textit{Physical Review D}, 97(10), 106017. \textit{arXiv:1802.10524}.

\bibitem{klitgaard2020}
Klitgaard N, Loll R. (2020).
How round is the quantum de Sitter universe?
\textit{European Physical Journal C}, 80, 990. \textit{arXiv:2006.06263}.

\bibitem{fortunato2007}
Fortunato S, Barthélemy M. (2007).
Resolution limit in community detection.
\textit{Proceedings of the National Academy of Sciences}, 104(1), 36--41.

\bibitem{chao2012coverage}
Chao A, Jost L. (2012).
Coverage-based rarefaction and extrapolation: standardizing samples by
completeness rather than size.
\textit{Ecology}, 93(12), 2533--2547.

\bibitem{clauset2009}
Clauset A, Shalizi C R, Newman M E J. (2009).
Power-law distributions in empirical data.
\textit{SIAM Review}, 51(4), 661--703.

\bibitem{heaps1978}
Heaps H S. (1978).
\textit{Information Retrieval: Computational and Theoretical Aspects}.
Academic Press. See also Herdan G. (1964), \textit{Quantitative Linguistics},
Butterworths.

\bibitem{gerlach2013}
Gerlach M, Altmann E G. (2013).
Stochastic model for the vocabulary growth in natural languages.
\textit{Physical Review X}, 3, 021006.

\bibitem{hubbell2001}
Hubbell S P. (2001).
\textit{The Unified Neutral Theory of Biodiversity and Biogeography}.
Princeton University Press.

\bibitem{saati2026}
Saati-Santamar\'{i}a Z, et al. (2026).
Soil Microbial Diversity and Network Organization Respond to Land Use and
Agricultural Inputs Worldwide. \textit{Global Change Biology}, 32, e70984.

\bibitem{wagg2019}
Wagg C, Schlaeppi K, Banerjee S, Kuramae E E, van der Heijden M G A. (2019).
Fungal-bacterial diversity and microbiome complexity predict ecosystem
functioning. \textit{Nature Communications}, 10, 4841.

\bibitem{herren2017}
Herren C M, McMahon K D. (2017).
Cohesion: a method for quantifying the connectivity of microbial communities.
\textit{The ISME Journal}, 11(11), 2426--2438.

\bibitem{yan2017}
Yan G, Martinez N D, Liu Y-Y. (2017).
Degree heterogeneity and stability of ecological networks.
\textit{Journal of the Royal Society Interface}, 14(131), 20170189.

\bibitem{yuan2021}
Yuan M M, Guo X, Wu L, et al. (2021).
Climate warming enhances microbial network complexity and stability.
\textit{Nature Climate Change}, 11, 343--348.

\bibitem{thebault2010}
Th\'{e}bault E, Fontaine C. (2010).
Stability of ecological communities and the architecture of mutualistic and
trophic networks. \textit{Science}, 329(5993), 853--856.

\bibitem{baldassano2016}
Baldassano S N, Bassett D S. (2016).
Topological distortion and reorganized modular structure of gut microbial
co-occurrence networks in inflammatory bowel disease.
\textit{Scientific Reports}, 6, 26087.

\bibitem{may1972}
May R M. (1972).
Will a large complex system be stable? \textit{Nature}, 238, 413--414.

\bibitem{allesina2012}
Allesina S, Tang S. (2012).
Stability criteria for complex ecosystems. \textit{Nature}, 483, 205--208.

\bibitem{morrien2017}
Morri\"{e}n E, et al. (2017).
Soil networks become more connected and take up more carbon as nature
restoration progresses. \textit{Nature Communications}, 8, 14349.

\bibitem{banerjee2019}
Banerjee S, et al. (2019).
Agricultural intensification reduces microbial network complexity and the
abundance of keystone taxa in roots. \textit{The ISME Journal}, 13, 1722--1736.

\bibitem{niu2026}
Niu R, Shu X, Zhao Y. (2026).
Reliability-Contagion Feasibility in LLM Multi-Agent Networks.
arXiv:2607.21912 [cs.MA].

\bibitem{ulanowicz2009}
Ulanowicz R E, Goerner S J, Lietaer B, Gomez R. (2009).
Quantifying sustainability: resilience, efficiency and the return of
information theory. \textit{Ecological Complexity}, 6(1), 27--36.

\bibitem{leskovec2005}
Leskovec J, Kleinberg J, Faloutsos C. (2005).
Graphs over time: densification laws, shrinking diameters and possible
explanations. \textit{Proceedings of the 11th ACM SIGKDD International
Conference on Knowledge Discovery and Data Mining}, 177--187.

\bibitem{gao2021}
Gao L, Rouz\'{e} C. (2021).
Ricci curvature of quantum channels on non-commutative transportation
metric spaces. \textit{arXiv:2108.10609}.

\end{thebibliography}

\appendix
% ============================================================
\section{The Two-Sided Reading Made Exact: Ascendency on an Unweighted Graph}
\label{app:ascendency}
% ============================================================

Remark~\ref{rem:window} records that Definition~\ref{def:fsri} is monotone in
each of its factors and therefore expresses one edge of a two-sided constraint
rather than the constraint itself, and names Ulanowicz's window of
vitality~\cite{ulanowicz2009} as the ecological antecedent. That framework is
stated for weighted flow networks and its transplant to the setting of this
paper is not obvious, so we carry it out here. The result is short, exact, and
it settles what such a transplant can and cannot deliver.

\subsection{Setup}

Let $A$ be the adjacency matrix of an undirected unweighted graph with $m$
edges and degree sequence $\mathbf{k} = A\mathbf{1}$. Ulanowicz's decomposition
takes a matrix of flows; on an unweighted graph the only choice that introduces
no free parameter is the uniform one,
\begin{equation}
p_{ij} \;=\; \frac{A_{ij}}{2m},
\qquad p_{i\cdot} \;=\; p_{\cdot i} \;=\; \frac{k_i}{2m}.
\label{eq:uniform_flow}
\end{equation}
The capacity is $H = -\sum_{ij} p_{ij}\log p_{ij}$, the ascendency density is
the average mutual information $\mathcal{A} = \sum_{ij} p_{ij}\log\!\big(p_{ij}
/ (p_{i\cdot}p_{\cdot j})\big)$, the overhead is $\Phi = H - \mathcal{A}$, and
the degree of order is $\alpha = \mathcal{A}/H$.

\subsection{Three exact statements}

\begin{lemma}[Capacity is fixed by the edge count]
\label{lem:capacity}
Under \eqref{eq:uniform_flow}, $H = \log(2m)$ for every graph.
\end{lemma}

\begin{proof}
There are exactly $2m$ non-zero entries of $p$, each equal to $1/2m$, so
\[
H = -2m \cdot \tfrac{1}{2m}\log\tfrac{1}{2m} = \log(2m). \qedhere
\]
\end{proof}

\begin{proposition}[Ascendency is mean surprisal against the configuration
null]
\label{prop:ascendency_null}
Let $P^0_{ij} = k_i k_j / 2m$ be the expected number of edges between $i$ and
$j$ under the configuration model. Then under \eqref{eq:uniform_flow}
\begin{equation}
\mathcal{A} \;=\; \frac{1}{2m}\sum_{(i,j):\,A_{ij}=1}
\log\frac{2m}{k_i k_j}
\;=\; \big\langle -\log P^0_{ij} \big\rangle_{A_{ij}=1},
\label{eq:ascendency_surprisal}
\end{equation}
the mean Shannon surprisal of the realised edges under that null, and
equivalently the Kullback-Leibler divergence of the edge distribution from the
product of its marginals.
\end{proposition}

\begin{proof}
Substituting \eqref{eq:uniform_flow} into the definition of $\mathcal{A}$, each
non-zero term contributes
\[
\frac{1}{2m}\log\!\left[\frac{1/2m}{k_ik_j/(2m)^2}\right]
\;=\; \frac{1}{2m}\log\!\left(\frac{2m}{k_ik_j}\right),
\]
and $2m/(k_ik_j) = 1/P^0_{ij}$.
\end{proof}

This identifies the construction rather than merely permitting it. The null
$P^0$ is the same one against which Newman modularity is defined; the two
quantities are different divergences from it, modularity taking the difference
$A_{ij} - P^0_{ij}$ and ascendency the logarithm of the ratio. The transplant
is therefore not an analogy imported from another field but a second reading of
an object already standard in network science.

\begin{corollary}[Closed form]
\label{cor:closed_form}
Write $\langle \log k\rangle_w = \sum_i (k_i/2m)\log k_i$ for the
degree-weighted mean of $\log k$. Then
\begin{equation}
\Phi = 2\langle \log k\rangle_w,
\qquad
\mathcal{A} = \log(2m) - 2\langle \log k\rangle_w,
\qquad
\alpha = 1 - \frac{2\langle \log k\rangle_w}{\log(2m)} .
\label{eq:closed_form}
\end{equation}
\end{corollary}

\begin{proof}
Expand \eqref{eq:ascendency_surprisal}: $\sum_{ij} A_{ij}\log k_i = \sum_i k_i
\log k_i$ and likewise for $j$, so the double sum contributes $2\sum_i k_i \log
k_i = 4m \langle \log k\rangle_w$. Divide by $2m$ and subtract from
$\log(2m)$; then use Lemma~\ref{lem:capacity}.
\end{proof}

\begin{corollary}[Optimal degree on regular graphs]
\label{cor:kstar}
For a $k$-regular graph on $n$ vertices, $\mathcal{A} = \log(n/k)$,
$\Phi = 2\log k$ and $\alpha = \log(n/k)/\log(nk)$. The robustness function
$R(\alpha) = -\alpha\log\alpha$ of~\cite{ulanowicz2009} is maximal at
$\alpha = 1/e$, which holds exactly when
\begin{equation}
\frac{\log k^\ast}{\log n} \;=\; \frac{e-1}{e+1} \;=\; 0.4621\ldots,
\qquad\text{that is}\qquad k^\ast = n^{(e-1)/(e+1)} .
\label{eq:kstar}
\end{equation}
\end{corollary}

Two features of \eqref{eq:kstar} are worth stating. The optimal degree grows
sublinearly with the order, so the optimal density $D^\ast = n^{-2/(e+1)}$
tends to zero: the two-sided reading prescribes increasing sparsity with size,
which is the regime in which large empirical networks are found. And
\eqref{eq:kstar} is a point prediction, falsifiable on any family of graphs
with varying order, which is more than the monotone reading of
Definition~\ref{def:fsri} can offer.

\subsection{What the transplant cannot do}

\begin{remark}[The order parameter is a degree-sequence statistic]
\label{rem:degree_only}
Equation~\eqref{eq:closed_form} depends on the graph only through $m$ and
$\mathbf{k}$. Hence $\mathcal{A}$, $\Phi$, $\alpha$ and $R$ are invariant under
any rewiring that preserves the degree sequence, and in particular they cannot
distinguish a clustered graph from its configuration-model randomisation. A
numerical check on double-edge swaps confirms this: the global clustering
coefficient moves while $\alpha$ is unchanged to six decimal places.

This bounds the reading of Remark~\ref{rem:window} more precisely than that
remark could. The window obtained by this route is a statement about
\emph{degree distributions} and not about wiring. It also corrects a natural
but wrong way to describe the pairing: $\alpha$ and $\mathrm{Coex}$ are not a
level coordinate and a shape coordinate drawn from different parts of the
graph, since both are functionals of $\mathbf{k}$ alone, differing in
weighting, logarithmic against quadratic, with $\mathrm{Coex}$ a pure
dispersion that vanishes on regular graphs and $\alpha$ moving with the level
as well. The genuine wiring content of Definition~\ref{def:fsri} is carried by
$C$ and $M$, and no quantity in this appendix touches it. A three-way reading
is therefore the accurate one: the degree sequence, read by $\alpha$ or by
$\mathrm{Coex}$; the wiring at fixed degrees, read by $C$ and $M$; and the
order.

Two further scope conditions apply. The uniform flow \eqref{eq:uniform_flow} is
a choice, and although it is the only parameter-free one available on an
unweighted graph, the quantities above are not the ascendency of a measured
flow network. And Ulanowicz's overhead decomposes into imports, exports,
dissipation and redundancy on a directed flow network, of which a symmetric $A$
retains only the last; $\Phi$ above should be read as a redundancy and not as
the full overhead.
\end{remark}

\end{document}
